% Môn: Vật lý
% Lớp: 10
% Chương: Chương III. Động lực học
% Bài: L10C3 Bài 22. Thực hành: Tổng hợp lực
% Số câu: 162
% App đọc file này trực tiếp — sửa rồi Commit trên GitHub, bấm Đồng bộ GitHub .tex

% L10C3 Bài 22. Thực hành: Tổng hợp lực

% ===== Câu 1 =====
% ID: AIQ_L10C3_FULL_1459
% Mức: NB
\dangbt{Nhận biết dụng cụ và bố trí thí nghiệm tổng hợp lực}
\begin{ex}
% ID: AIQ_L10C3_FULL_1459
% Mức: NB
Dụng cụ dùng để đo độ lớn của lực trong thí nghiệm tổng hợp lực là gì?
\choice
{\True Lực kế}
{Thước đo góc}
{Đồng hồ bấm giây}
{Nhiệt kế}
\loigiai{
Lực kế cho số chỉ độ lớn của lực.
}
\end{ex}

% ===== Câu 2 =====
% ID: AIQ_L10C3_FULL_1460
% Mức: NB
\begin{ex}
% ID: AIQ_L10C3_FULL_1460
% Mức: NB
Tấm bảng trong thí nghiệm nên được bố trí như thế nào?
\choice
{Nằm tự do trên bàn}
{\True Thẳng đứng và cố định}
{Nghiêng và rung được}
{Treo bằng một sợi chỉ}
\loigiai{
Bảng thẳng đứng, cố định giúp xác định phương lực ổn định.
}
\end{ex}

% ===== Câu 3 =====
% ID: AIQ_L10C3_FULL_1461
% Mức: NB
\begin{ex}
% ID: AIQ_L10C3_FULL_1461
% Mức: NB
Vòng nhẫn nhỏ ở nút dây có vai trò chính gì?
\choice
{Đo góc giữa hai lực}
{Thay thế lực kế}
{\True Điểm đồng quy của các lực}
{Giảm khối lượng vật}
\loigiai{
Các dây cùng tác dụng vào vòng nhẫn nên các lực đồng quy tại đó.
}
\end{ex}

% ===== Câu 4 =====
% ID: AIQ_L10C3_FULL_1462
% Mức: TH
\begin{ex}
% ID: AIQ_L10C3_FULL_1462
% Mức: TH
Trước khi đo bằng lực kế cần làm gì?
\choice
{Kéo lực kế quá giới hạn}
{Đặt lực kế nằm ngoài mặt phẳng dây}
{Tháo lò xo khỏi lực kế}
{\True Kiểm tra và chỉnh vạch số 0}
\loigiai{
Chỉnh số 0 làm giảm sai số hệ thống.
}
\end{ex}

% ===== Câu 5 =====
% ID: AIQ_L10C3_FULL_1463
% Mức: TH
\begin{ex}
% ID: AIQ_L10C3_FULL_1463
% Mức: TH
Khi đọc lực kế, mắt phải đặt thế nào?
\choice
{\True Vuông góc với mặt chia độ}
{Nhìn xiên từ bên trái}
{Nhìn xiên từ bên phải}
{Đặt sát vào lò xo}
\loigiai{
Đặt mắt vuông góc giúp tránh sai số thị sai.
}
\end{ex}

% ===== Câu 6 =====
% ID: AIQ_L10C3_FULL_1464
% Mức: TH
\begin{ex}
% ID: AIQ_L10C3_FULL_1464
% Mức: TH
Các dây nối vòng nhẫn với lực kế phải có đặc điểm nào?
\choice
{Chùng và khác mặt phẳng}
{\True Căng và cùng nằm trong một mặt phẳng}
{Quấn nhiều vòng quanh móc}
{Có độ dài bằng nhau tuyệt đối}
\loigiai{
Dây căng, đồng phẳng bảo đảm phương lực rõ ràng.
}
\end{ex}

% ===== Câu 7 =====
% ID: AIQ_L10C3_FULL_1465
% Mức: VD
\begin{ex}
% ID: AIQ_L10C3_FULL_1465
% Mức: VD
Giới hạn đo của lực kế phải được chọn thế nào?
\choice
{Nhỏ hơn lực dự kiến}
{Bằng 0}
{\True Lớn hơn giá trị lực dự kiến}
{Không cần quan tâm}
\loigiai{
Chọn giới hạn đo phù hợp để tránh quá tải.
}
\end{ex}

% ===== Câu 8 =====
% ID: AIQ_L10C3_FULL_1466
% Mức: VD
\begin{ex}
% ID: AIQ_L10C3_FULL_1466
% Mức: VD
Để đánh dấu phương của lực trên giấy, cần dựa vào yếu tố nào?
\choice
{Màu của lực kế}
{Chiều dài bảng}
{Khối lượng vòng nhẫn}
{\True Phương của đoạn dây căng}
\loigiai{
Lực căng tác dụng dọc theo dây.
}
\end{ex}

% ===== Câu 9 =====
% ID: AIQ_L10C3_FULL_1467
% Mức: VD
\begin{ex}
% ID: AIQ_L10C3_FULL_1467
% Mức: VD
Trạng thái cân bằng được nhận biết khi nào?
\choice
{\True Vòng nhẫn đứng yên tại vị trí đánh dấu}
{Một lực kế chỉ 0}
{Các dây đều chùng}
{Bảng bắt đầu rung}
\loigiai{
Vòng nhẫn đứng yên tại mốc cho thấy hệ lực cân bằng.
}
\end{ex}

% ===== Câu 10 =====
% ID: AIQ_L10C3_FULL_1468
% Mức: VDC
\begin{ex}
% ID: AIQ_L10C3_FULL_1468
% Mức: VDC
Đơn vị thường dùng trên lực kế trong thí nghiệm là gì?
\choice
{kg}
{\True $N$}
{m}
{s}
\loigiai{
Đơn vị SI của lực là niutơn, kí hiệu N.
}
\end{ex}

% ===== Câu 11 =====
% ID: AIQ_L10C3_FULL_1469
% Mức: NB
\begin{ex}
% ID: AIQ_L10C3_FULL_1469
% Mức: NB
Xét bố trí thí nghiệm tổng hợp lực lần 1.
\choiceTF
{\True Lực kế dùng để đo độ lớn lực.}
{\True Dây căng cho biết phương của lực.}
{Có thể đọc lực kế khi mắt nhìn xiên.}
{\True Vòng nhẫn là điểm đồng quy của các lực.}
\loigiai{
\begin{itemchoice}
\item (Đúng) Lực kế dùng để đo độ lớn lực. (Vì): lực kế đo độ lớn lực. b) Đúng: lực căng có phương dọc theo dây. c) Sai: nhìn xiên gây sai số thị sai. d) Đúng: các lực đồng quy tại vòng nhẫn
\item (Đúng) Dây căng cho biết phương của lực.
\item (Sai) Có thể đọc lực kế khi mắt nhìn xiên.
\item (Đúng) Vòng nhẫn là điểm đồng quy của các lực.
\end{itemchoice}
}
\end{ex}

% ===== Câu 12 =====
% ID: AIQ_L10C3_FULL_1470
% Mức: TH
\begin{ex}
% ID: AIQ_L10C3_FULL_1470
% Mức: TH
Xét bố trí thí nghiệm tổng hợp lực lần 2.
\choiceTF
{\True Lực kế dùng để đo độ lớn lực.}
{\True Dây căng cho biết phương của lực.}
{Có thể đọc lực kế khi mắt nhìn xiên.}
{\True Vòng nhẫn là điểm đồng quy của các lực.}
\loigiai{
\begin{itemchoice}
\item (Đúng) Lực kế dùng để đo độ lớn lực. (Vì): lực kế đo độ lớn lực. b) Đúng: lực căng có phương dọc theo dây. c) Sai: nhìn xiên gây sai số thị sai. d) Đúng: các lực đồng quy tại vòng nhẫn
\item (Đúng) Dây căng cho biết phương của lực.
\item (Sai) Có thể đọc lực kế khi mắt nhìn xiên.
\item (Đúng) Vòng nhẫn là điểm đồng quy của các lực.
\end{itemchoice}
}
\end{ex}

% ===== Câu 13 =====
% ID: AIQ_L10C3_FULL_1471
% Mức: VD
\begin{ex}
% ID: AIQ_L10C3_FULL_1471
% Mức: VD
Xét bố trí thí nghiệm tổng hợp lực lần 3.
\choiceTF
{\True Lực kế dùng để đo độ lớn lực.}
{\True Dây căng cho biết phương của lực.}
{Có thể đọc lực kế khi mắt nhìn xiên.}
{\True Vòng nhẫn là điểm đồng quy của các lực.}
\loigiai{
\begin{itemchoice}
\item (Đúng) Lực kế dùng để đo độ lớn lực. (Vì): lực kế đo độ lớn lực. b) Đúng: lực căng có phương dọc theo dây. c) Sai: nhìn xiên gây sai số thị sai. d) Đúng: các lực đồng quy tại vòng nhẫn
\item (Đúng) Dây căng cho biết phương của lực.
\item (Sai) Có thể đọc lực kế khi mắt nhìn xiên.
\item (Đúng) Vòng nhẫn là điểm đồng quy của các lực.
\end{itemchoice}
}
\end{ex}

% ===== Câu 14 =====
% ID: AIQ_L10C3_FULL_1472
% Mức: VD
\begin{ex}
% ID: AIQ_L10C3_FULL_1472
% Mức: VD
Xét bố trí thí nghiệm tổng hợp lực lần 4.
\choiceTF
{\True Lực kế dùng để đo độ lớn lực.}
{\True Dây căng cho biết phương của lực.}
{Có thể đọc lực kế khi mắt nhìn xiên.}
{\True Vòng nhẫn là điểm đồng quy của các lực.}
\loigiai{
\begin{itemchoice}
\item (Đúng) Lực kế dùng để đo độ lớn lực. (Vì): lực kế đo độ lớn lực. b) Đúng: lực căng có phương dọc theo dây. c) Sai: nhìn xiên gây sai số thị sai. d) Đúng: các lực đồng quy tại vòng nhẫn
\item (Đúng) Dây căng cho biết phương của lực.
\item (Sai) Có thể đọc lực kế khi mắt nhìn xiên.
\item (Đúng) Vòng nhẫn là điểm đồng quy của các lực.
\end{itemchoice}
}
\end{ex}

% ===== Câu 15 =====
% ID: AIQ_L10C3_FULL_1473
% Mức: VDC
\begin{ex}
% ID: AIQ_L10C3_FULL_1473
% Mức: VDC
Xét bố trí thí nghiệm tổng hợp lực lần 5.
\choiceTF
{\True Lực kế dùng để đo độ lớn lực.}
{\True Dây căng cho biết phương của lực.}
{Có thể đọc lực kế khi mắt nhìn xiên.}
{\True Vòng nhẫn là điểm đồng quy của các lực.}
\loigiai{
\begin{itemchoice}
\item (Đúng) Lực kế dùng để đo độ lớn lực. (Vì): lực kế đo độ lớn lực. b) Đúng: lực căng có phương dọc theo dây. c) Sai: nhìn xiên gây sai số thị sai. d) Đúng: các lực đồng quy tại vòng nhẫn
\item (Đúng) Dây căng cho biết phương của lực.
\item (Sai) Có thể đọc lực kế khi mắt nhìn xiên.
\item (Đúng) Vòng nhẫn là điểm đồng quy của các lực.
\end{itemchoice}
}
\end{ex}

% ===== Câu 16 =====
% ID: AIQ_L10C3_FULL_1474
% Mức: TH
\begin{ex}
% ID: AIQ_L10C3_FULL_1474
% Mức: TH
Trong một bố trí có 3 lực kế nối với cùng một vòng nhẫn. Có bao nhiêu lực do các lực kế tác dụng trực tiếp lên vòng nhẫn?
\shortans{3}
\loigiai{
Mỗi lực kế truyền một lực căng qua một dây đến vòng nhẫn. Vì vậy số lực tác dụng trực tiếp bằng số lực kế, kết quả là 3. Đáp án không ghi đơn vị.
}
\end{ex}

% ===== Câu 17 =====
% ID: AIQ_L10C3_FULL_1475
% Mức: TH
\begin{ex}
% ID: AIQ_L10C3_FULL_1475
% Mức: TH
Trong một bố trí có 4 lực kế nối với cùng một vòng nhẫn. Có bao nhiêu lực do các lực kế tác dụng trực tiếp lên vòng nhẫn?
\shortans{4}
\loigiai{
Mỗi lực kế truyền một lực căng qua một dây đến vòng nhẫn. Vì vậy số lực tác dụng trực tiếp bằng số lực kế, kết quả là 4. Đáp án không ghi đơn vị.
}
\end{ex}

% ===== Câu 18 =====
% ID: AIQ_L10C3_FULL_1476
% Mức: VD
\begin{ex}
% ID: AIQ_L10C3_FULL_1476
% Mức: VD
Trong một bố trí có 5 lực kế nối với cùng một vòng nhẫn. Có bao nhiêu lực do các lực kế tác dụng trực tiếp lên vòng nhẫn?
\shortans{5}
\loigiai{
Mỗi lực kế truyền một lực căng qua một dây đến vòng nhẫn. Vì vậy số lực tác dụng trực tiếp bằng số lực kế, kết quả là 5. Đáp án không ghi đơn vị.
}
\end{ex}

% ===== Câu 19 =====
% ID: AIQ_L10C3_FULL_1477
% Mức: VD
\begin{ex}
% ID: AIQ_L10C3_FULL_1477
% Mức: VD
Trong một bố trí có 6 lực kế nối với cùng một vòng nhẫn. Có bao nhiêu lực do các lực kế tác dụng trực tiếp lên vòng nhẫn?
\shortans{6}
\loigiai{
Mỗi lực kế truyền một lực căng qua một dây đến vòng nhẫn. Vì vậy số lực tác dụng trực tiếp bằng số lực kế, kết quả là 6. Đáp án không ghi đơn vị.
}
\end{ex}

% ===== Câu 20 =====
% ID: AIQ_L10C3_FULL_1478
% Mức: VD
\begin{ex}
% ID: AIQ_L10C3_FULL_1478
% Mức: VD
Trong một bố trí có 7 lực kế nối với cùng một vòng nhẫn. Có bao nhiêu lực do các lực kế tác dụng trực tiếp lên vòng nhẫn?
\shortans{7}
\loigiai{
Mỗi lực kế truyền một lực căng qua một dây đến vòng nhẫn. Vì vậy số lực tác dụng trực tiếp bằng số lực kế, kết quả là 7. Đáp án không ghi đơn vị.
}
\end{ex}

% ===== Câu 21 =====
% ID: AIQ_L10C3_FULL_1479
% Mức: VDC
\begin{ex}
% ID: AIQ_L10C3_FULL_1479
% Mức: VDC
Trong một bố trí có 8 lực kế nối với cùng một vòng nhẫn. Có bao nhiêu lực do các lực kế tác dụng trực tiếp lên vòng nhẫn?
\shortans{8}
\loigiai{
Mỗi lực kế truyền một lực căng qua một dây đến vòng nhẫn. Vì vậy số lực tác dụng trực tiếp bằng số lực kế, kết quả là 8. Đáp án không ghi đơn vị.
}
\end{ex}

% ===== Câu 22 =====
% ID: AIQ_L10C3_FULL_1480
% Mức: TH
\begin{ex}
% ID: AIQ_L10C3_FULL_1480
% Mức: TH
Trình bày cách bố trí dụng cụ để khảo sát tổng hợp hai lực đồng quy (phương án 1).
\choiceTF

\loigiai{
Cố định bảng; kiểm tra và chỉnh số 0 các lực kế; nối các dây vào cùng vòng nhẫn; bố trí dây căng, đồng phẳng; đánh dấu vị trí cân bằng và phương từng dây; đọc lực kế với mắt vuông góc mặt chia.
}
\end{ex}

% ===== Câu 23 =====
% ID: AIQ_L10C3_FULL_1481
% Mức: TH
\begin{ex}
% ID: AIQ_L10C3_FULL_1481
% Mức: TH
Trình bày cách bố trí dụng cụ để khảo sát tổng hợp hai lực đồng quy (phương án 2).
\choiceTF

\loigiai{
Cố định bảng; kiểm tra và chỉnh số 0 các lực kế; nối các dây vào cùng vòng nhẫn; bố trí dây căng, đồng phẳng; đánh dấu vị trí cân bằng và phương từng dây; đọc lực kế với mắt vuông góc mặt chia.
}
\end{ex}

% ===== Câu 24 =====
% ID: AIQ_L10C3_FULL_1482
% Mức: VD
\begin{ex}
% ID: AIQ_L10C3_FULL_1482
% Mức: VD
Trình bày cách bố trí dụng cụ để khảo sát tổng hợp hai lực đồng quy (phương án 3).
\choiceTF

\loigiai{
Cố định bảng; kiểm tra và chỉnh số 0 các lực kế; nối các dây vào cùng vòng nhẫn; bố trí dây căng, đồng phẳng; đánh dấu vị trí cân bằng và phương từng dây; đọc lực kế với mắt vuông góc mặt chia.
}
\end{ex}

% ===== Câu 25 =====
% ID: AIQ_L10C3_FULL_1483
% Mức: VD
\begin{ex}
% ID: AIQ_L10C3_FULL_1483
% Mức: VD
Trình bày cách bố trí dụng cụ để khảo sát tổng hợp hai lực đồng quy (phương án 4).
\choiceTF

\loigiai{
Cố định bảng; kiểm tra và chỉnh số 0 các lực kế; nối các dây vào cùng vòng nhẫn; bố trí dây căng, đồng phẳng; đánh dấu vị trí cân bằng và phương từng dây; đọc lực kế với mắt vuông góc mặt chia.
}
\end{ex}

% ===== Câu 26 =====
% ID: AIQ_L10C3_FULL_1484
% Mức: VD
\begin{ex}
% ID: AIQ_L10C3_FULL_1484
% Mức: VD
Trình bày cách bố trí dụng cụ để khảo sát tổng hợp hai lực đồng quy (phương án 5).
\choiceTF

\loigiai{
Cố định bảng; kiểm tra và chỉnh số 0 các lực kế; nối các dây vào cùng vòng nhẫn; bố trí dây căng, đồng phẳng; đánh dấu vị trí cân bằng và phương từng dây; đọc lực kế với mắt vuông góc mặt chia.
}
\end{ex}

% ===== Câu 27 =====
% ID: AIQ_L10C3_FULL_1485
% Mức: VDC
\begin{ex}
% ID: AIQ_L10C3_FULL_1485
% Mức: VDC
Trình bày cách bố trí dụng cụ để khảo sát tổng hợp hai lực đồng quy (phương án 6).
\choiceTF

\loigiai{
Cố định bảng; kiểm tra và chỉnh số 0 các lực kế; nối các dây vào cùng vòng nhẫn; bố trí dây căng, đồng phẳng; đánh dấu vị trí cân bằng và phương từng dây; đọc lực kế với mắt vuông góc mặt chia.
}
\end{ex}

% ===== Câu 28 =====
% ID: AIQ_L10C3_FULL_1486
% Mức: NB
\dangbt{Đọc và xử lí số liệu đo lực}
\begin{ex}
% ID: AIQ_L10C3_FULL_1486
% Mức: NB
Một lực kế có độ chia nhỏ nhất $0,1\ \mathrm{N}$. Cách ghi nào phù hợp?
\choice
{\True $2,3\ \mathrm{N}$}
{$2,3456\ \mathrm{N}$}
{$23\ \mathrm{N}$}
{$0,23\ \mathrm{N}$}
\loigiai{
Kết quả được ghi phù hợp với độ chia nhỏ nhất của dụng cụ.
}
\end{ex}

% ===== Câu 29 =====
% ID: AIQ_L10C3_FULL_1487
% Mức: NB
\begin{ex}
% ID: AIQ_L10C3_FULL_1487
% Mức: NB
Ba lần đo một lực cho $2{,}1N$; $2{,}3N$; $2{,}2$ N. Giá trị trung bình là bao nhiêu?
\choice
{$2,0\ \mathrm{N}$}
{\True $2,2\ \mathrm{N}$}
{$2,3\ \mathrm{N}$}
{$6,6\ \mathrm{N}$}
\loigiai{
$\overline F=(2,1+2,3+2,2)/3=2,2\ \mathrm{N}$.
}
\end{ex}

% ===== Câu 30 =====
% ID: AIQ_L10C3_FULL_1488
% Mức: NB
\begin{ex}
% ID: AIQ_L10C3_FULL_1488
% Mức: NB
Khi số chỉ lực kế dao động nhẹ, nên xử lí thế nào?
\choice
{Chọn ngay giá trị lớn nhất}
{Kéo mạnh thêm}
{\True Chờ ổn định rồi đọc nhiều lần}
{Ghi bằng 0}
\loigiai{
Đọc lặp lại khi hệ ổn định làm kết quả tin cậy hơn.
}
\end{ex}

% ===== Câu 31 =====
% ID: AIQ_L10C3_FULL_1489
% Mức: TH
\begin{ex}
% ID: AIQ_L10C3_FULL_1489
% Mức: TH
Số liệu $1{,}8N$; $1{,}9N$; $5{,}2N$ có một giá trị bất thường. Giá trị đó là gì?
\choice
{$1,8\ \mathrm{N}$}
{$1,9\ \mathrm{N}$}
{Không có}
{\True $5,2\ \mathrm{N}$}
\loigiai{
$5{,}2N$ lệch xa hai giá trị còn lại nên cần kiểm tra lại.
}
\end{ex}

% ===== Câu 32 =====
% ID: AIQ_L10C3_FULL_1490
% Mức: TH
\begin{ex}
% ID: AIQ_L10C3_FULL_1490
% Mức: TH
Góc giữa hai phương lực được đo bằng dụng cụ nào?
\choice
{\True Thước đo góc}
{Cân điện tử}
{Ampe kế}
{Nhiệt kế}
\loigiai{
Thước đo góc dùng để xác định góc giữa các phương lực.
}
\end{ex}

% ===== Câu 33 =====
% ID: AIQ_L10C3_FULL_1491
% Mức: TH
\begin{ex}
% ID: AIQ_L10C3_FULL_1491
% Mức: TH
Nếu lực kế chỉ giữa vạch $2{,}4N$ và $2{,}5N$, độ chia nhỏ nhất $0{,}1N$, nên ghi gần đúng bao nhiêu?
\choice
{$2,0\ \mathrm{N}$}
{\True $2,45\ \mathrm{N}$}
{$2,9\ \mathrm{N}$}
{$24,5\ \mathrm{N}$}
\loigiai{
Có thể ước lượng thêm một chữ số so với vạch chia.
}
\end{ex}

% ===== Câu 34 =====
% ID: AIQ_L10C3_FULL_1492
% Mức: VD
\begin{ex}
% ID: AIQ_L10C3_FULL_1492
% Mức: VD
Đại lượng nào cần ghi kèm khi đo lực?
\choice
{Tên người đo}
{Màu dây}
{\True Đơn vị đo}
{Kích thước phòng}
\loigiai{
Số đo vật lí phải có giá trị và đơn vị.
}
\end{ex}

% ===== Câu 35 =====
% ID: AIQ_L10C3_FULL_1493
% Mức: VD
\begin{ex}
% ID: AIQ_L10C3_FULL_1493
% Mức: VD
Khi lặp phép đo, mục đích chính là gì?
\choice
{Làm tăng giới hạn đo}
{Đổi đơn vị lực}
{Thay đổi phương lực}
{\True Giảm ảnh hưởng của sai số ngẫu nhiên}
\loigiai{
Lấy nhiều lần và tính trung bình giảm ảnh hưởng dao động ngẫu nhiên.
}
\end{ex}

% ===== Câu 36 =====
% ID: AIQ_L10C3_FULL_1494
% Mức: VD
\begin{ex}
% ID: AIQ_L10C3_FULL_1494
% Mức: VD
Bảng số liệu hợp lí cần có nội dung nào?
\choice
{\True Tên đại lượng, giá trị và đơn vị}
{Chỉ có số thứ tự}
{Chỉ có đơn vị}
{Chỉ có nhận xét}
\loigiai{
Tiêu đề đại lượng và đơn vị giúp dữ liệu rõ nghĩa.
}
\end{ex}

% ===== Câu 37 =====
% ID: AIQ_L10C3_FULL_1495
% Mức: VDC
\begin{ex}
% ID: AIQ_L10C3_FULL_1495
% Mức: VDC
Hai lần đo góc được $59^\circ$ và $61^\circ$. Giá trị trung bình là bao nhiêu?
\choice
{$59^\circ$}
{\True $60^\circ$}
{$61^\circ$}
{$120^\circ$}
\loigiai{
$\overline\alpha=(59+61)/2=60^\circ$.
}
\end{ex}

% ===== Câu 38 =====
% ID: AIQ_L10C3_FULL_1496
% Mức: NB
\begin{ex}
% ID: AIQ_L10C3_FULL_1496
% Mức: NB
Đánh giá cách đọc và xử lí số liệu ở lượt đo 1.
\choiceTF
{\True Số đo phải ghi kèm đơn vị.}
{\True Lấy trung bình nhiều lần giúp giảm sai số ngẫu nhiên.}
{Giá trị lệch xa các lần còn lại luôn được giữ mà không kiểm tra.}
{\True Độ chia nhỏ nhất ảnh hưởng cách ghi kết quả.}
\loigiai{
\begin{itemchoice}
\item (Đúng) Số đo phải ghi kèm đơn vị. (Vì): số đo phải kèm đơn vị. b) Đúng: lấy trung bình làm giảm ảnh hưởng sai số ngẫu nhiên. c) Sai: giá trị bất thường phải được kiểm tra và đo lại. d) Đúng: độ chia nhỏ nhất quyết định độ phân giải
\item (Đúng) Lấy trung bình nhiều lần giúp giảm sai số ngẫu nhiên.
\item (Sai) Giá trị lệch xa các lần còn lại luôn được giữ mà không kiểm tra.
\item (Đúng) Độ chia nhỏ nhất ảnh hưởng cách ghi kết quả.
\end{itemchoice}
}
\end{ex}

% ===== Câu 39 =====
% ID: AIQ_L10C3_FULL_1497
% Mức: TH
\begin{ex}
% ID: AIQ_L10C3_FULL_1497
% Mức: TH
Đánh giá cách đọc và xử lí số liệu ở lượt đo 2.
\choiceTF
{\True Số đo phải ghi kèm đơn vị.}
{\True Lấy trung bình nhiều lần giúp giảm sai số ngẫu nhiên.}
{Giá trị lệch xa các lần còn lại luôn được giữ mà không kiểm tra.}
{\True Độ chia nhỏ nhất ảnh hưởng cách ghi kết quả.}
\loigiai{
\begin{itemchoice}
\item (Đúng) Số đo phải ghi kèm đơn vị. (Vì): số đo phải kèm đơn vị. b) Đúng: lấy trung bình làm giảm ảnh hưởng sai số ngẫu nhiên. c) Sai: giá trị bất thường phải được kiểm tra và đo lại. d) Đúng: độ chia nhỏ nhất quyết định độ phân giải
\item (Đúng) Lấy trung bình nhiều lần giúp giảm sai số ngẫu nhiên.
\item (Sai) Giá trị lệch xa các lần còn lại luôn được giữ mà không kiểm tra.
\item (Đúng) Độ chia nhỏ nhất ảnh hưởng cách ghi kết quả.
\end{itemchoice}
}
\end{ex}

% ===== Câu 40 =====
% ID: AIQ_L10C3_FULL_1498
% Mức: VD
\begin{ex}
% ID: AIQ_L10C3_FULL_1498
% Mức: VD
Đánh giá cách đọc và xử lí số liệu ở lượt đo 3.
\choiceTF
{\True Số đo phải ghi kèm đơn vị.}
{\True Lấy trung bình nhiều lần giúp giảm sai số ngẫu nhiên.}
{Giá trị lệch xa các lần còn lại luôn được giữ mà không kiểm tra.}
{\True Độ chia nhỏ nhất ảnh hưởng cách ghi kết quả.}
\loigiai{
\begin{itemchoice}
\item (Đúng) Số đo phải ghi kèm đơn vị. (Vì): số đo phải kèm đơn vị. b) Đúng: lấy trung bình làm giảm ảnh hưởng sai số ngẫu nhiên. c) Sai: giá trị bất thường phải được kiểm tra và đo lại. d) Đúng: độ chia nhỏ nhất quyết định độ phân giải
\item (Đúng) Lấy trung bình nhiều lần giúp giảm sai số ngẫu nhiên.
\item (Sai) Giá trị lệch xa các lần còn lại luôn được giữ mà không kiểm tra.
\item (Đúng) Độ chia nhỏ nhất ảnh hưởng cách ghi kết quả.
\end{itemchoice}
}
\end{ex}

% ===== Câu 41 =====
% ID: AIQ_L10C3_FULL_1499
% Mức: VD
\begin{ex}
% ID: AIQ_L10C3_FULL_1499
% Mức: VD
Đánh giá cách đọc và xử lí số liệu ở lượt đo 4.
\choiceTF
{\True Số đo phải ghi kèm đơn vị.}
{\True Lấy trung bình nhiều lần giúp giảm sai số ngẫu nhiên.}
{Giá trị lệch xa các lần còn lại luôn được giữ mà không kiểm tra.}
{\True Độ chia nhỏ nhất ảnh hưởng cách ghi kết quả.}
\loigiai{
\begin{itemchoice}
\item (Đúng) Số đo phải ghi kèm đơn vị. (Vì): số đo phải kèm đơn vị. b) Đúng: lấy trung bình làm giảm ảnh hưởng sai số ngẫu nhiên. c) Sai: giá trị bất thường phải được kiểm tra và đo lại. d) Đúng: độ chia nhỏ nhất quyết định độ phân giải
\item (Đúng) Lấy trung bình nhiều lần giúp giảm sai số ngẫu nhiên.
\item (Sai) Giá trị lệch xa các lần còn lại luôn được giữ mà không kiểm tra.
\item (Đúng) Độ chia nhỏ nhất ảnh hưởng cách ghi kết quả.
\end{itemchoice}
}
\end{ex}

% ===== Câu 42 =====
% ID: AIQ_L10C3_FULL_1500
% Mức: VDC
\begin{ex}
% ID: AIQ_L10C3_FULL_1500
% Mức: VDC
Đánh giá cách đọc và xử lí số liệu ở lượt đo 5.
\choiceTF
{\True Số đo phải ghi kèm đơn vị.}
{\True Lấy trung bình nhiều lần giúp giảm sai số ngẫu nhiên.}
{Giá trị lệch xa các lần còn lại luôn được giữ mà không kiểm tra.}
{\True Độ chia nhỏ nhất ảnh hưởng cách ghi kết quả.}
\loigiai{
\begin{itemchoice}
\item (Đúng) Số đo phải ghi kèm đơn vị. (Vì): số đo phải kèm đơn vị. b) Đúng: lấy trung bình làm giảm ảnh hưởng sai số ngẫu nhiên. c) Sai: giá trị bất thường phải được kiểm tra và đo lại. d) Đúng: độ chia nhỏ nhất quyết định độ phân giải
\item (Đúng) Lấy trung bình nhiều lần giúp giảm sai số ngẫu nhiên.
\item (Sai) Giá trị lệch xa các lần còn lại luôn được giữ mà không kiểm tra.
\item (Đúng) Độ chia nhỏ nhất ảnh hưởng cách ghi kết quả.
\end{itemchoice}
}
\end{ex}

% ===== Câu 43 =====
% ID: AIQ_L10C3_FULL_1501
% Mức: TH
\begin{ex}
% ID: AIQ_L10C3_FULL_1501
% Mức: TH
Ba lần đo cùng một lực được $2N$, $2,1N$ và $2,2$ N. Tính giá trị trung bình theo N.
\shortans{2,1}
\loigiai{
Cộng ba số đo rồi chia cho 3: $\overline F=(F_1+F_2+F_3)/3=2,1\ \mathrm{N}$. Đáp án không ghi đơn vị.
}
\end{ex}

% ===== Câu 44 =====
% ID: AIQ_L10C3_FULL_1502
% Mức: TH
\begin{ex}
% ID: AIQ_L10C3_FULL_1502
% Mức: TH
Ba lần đo cùng một lực được $2,2N$, $2,3N$ và $2,4$ N. Tính giá trị trung bình theo N.
\shortans{2,3}
\loigiai{
Cộng ba số đo rồi chia cho 3: $\overline F=(F_1+F_2+F_3)/3=2,3\ \mathrm{N}$. Đáp án không ghi đơn vị.
}
\end{ex}

% ===== Câu 45 =====
% ID: AIQ_L10C3_FULL_1503
% Mức: VD
\begin{ex}
% ID: AIQ_L10C3_FULL_1503
% Mức: VD
Ba lần đo cùng một lực được $2,4N$, $2,5N$ và $2,6$ N. Tính giá trị trung bình theo N.
\shortans{2,5}
\loigiai{
Cộng ba số đo rồi chia cho 3: $\overline F=(F_1+F_2+F_3)/3=2,5\ \mathrm{N}$. Đáp án không ghi đơn vị.
}
\end{ex}

% ===== Câu 46 =====
% ID: AIQ_L10C3_FULL_1504
% Mức: VD
\begin{ex}
% ID: AIQ_L10C3_FULL_1504
% Mức: VD
Ba lần đo cùng một lực được $2,6N$, $2,7N$ và $2,8$ N. Tính giá trị trung bình theo N.
\shortans{2,7}
\loigiai{
Cộng ba số đo rồi chia cho 3: $\overline F=(F_1+F_2+F_3)/3=2,7\ \mathrm{N}$. Đáp án không ghi đơn vị.
}
\end{ex}

% ===== Câu 47 =====
% ID: AIQ_L10C3_FULL_1505
% Mức: VD
\begin{ex}
% ID: AIQ_L10C3_FULL_1505
% Mức: VD
Ba lần đo cùng một lực được $2,8N$, $2,9N$ và $3$ N. Tính giá trị trung bình theo N.
\shortans{2,9}
\loigiai{
Cộng ba số đo rồi chia cho 3: $\overline F=(F_1+F_2+F_3)/3=2,9\ \mathrm{N}$. Đáp án không ghi đơn vị.
}
\end{ex}

% ===== Câu 48 =====
% ID: AIQ_L10C3_FULL_1506
% Mức: VDC
\begin{ex}
% ID: AIQ_L10C3_FULL_1506
% Mức: VDC
Ba lần đo cùng một lực được $3N$, $3,1N$ và $3,2$ N. Tính giá trị trung bình theo N.
\shortans{3,1}
\loigiai{
Cộng ba số đo rồi chia cho 3: $\overline F=(F_1+F_2+F_3)/3=3,1\ \mathrm{N}$. Đáp án không ghi đơn vị.
}
\end{ex}

% ===== Câu 49 =====
% ID: AIQ_L10C3_FULL_1507
% Mức: TH
\begin{bt}
% ID: AIQ_L10C3_FULL_1507
% Mức: TH
Ba lần đo lực được $2{,}1N$, $2{,}2N$ và $2{,}3$ N. Hãy tính trung bình và nêu cách ghi kết quả.
\loigiai{
Tính trung bình cộng của ba số đo theo biểu thức $\overline F=(F_1+F_2+F_3)/3$. Kết quả là 2,2 $N$; khi ghi phải kèm đơn vị và làm tròn phù hợp độ chia nhỏ nhất.
}
\end{bt}

% ===== Câu 50 =====
% ID: AIQ_L10C3_FULL_1508
% Mức: TH
\begin{bt}
% ID: AIQ_L10C3_FULL_1508
% Mức: TH
Ba lần đo lực được $2{,}2N$, $2{,}3N$ và $2{,}4$ N. Hãy tính trung bình và nêu cách ghi kết quả.
\loigiai{
Tính trung bình cộng của ba số đo theo biểu thức $\overline F=(F_1+F_2+F_3)/3$. Kết quả là 2,3 $N$; khi ghi phải kèm đơn vị và làm tròn phù hợp độ chia nhỏ nhất.
}
\end{bt}

% ===== Câu 51 =====
% ID: AIQ_L10C3_FULL_1509
% Mức: VD
\begin{bt}
% ID: AIQ_L10C3_FULL_1509
% Mức: VD
Ba lần đo lực được $2{,}3N$, $2{,}4N$ và $2{,}5$ N. Hãy tính trung bình và nêu cách ghi kết quả.
\loigiai{
Tính trung bình cộng của ba số đo theo biểu thức $\overline F=(F_1+F_2+F_3)/3$. Kết quả là 2,4 $N$; khi ghi phải kèm đơn vị và làm tròn phù hợp độ chia nhỏ nhất.
}
\end{bt}

% ===== Câu 52 =====
% ID: AIQ_L10C3_FULL_1510
% Mức: VD
\begin{bt}
% ID: AIQ_L10C3_FULL_1510
% Mức: VD
Ba lần đo lực được $2{,}4N$, $2{,}5N$ và $2{,}6$ N. Hãy tính trung bình và nêu cách ghi kết quả.
\loigiai{
Tính trung bình cộng của ba số đo theo biểu thức $\overline F=(F_1+F_2+F_3)/3$. Kết quả là 2,5 $N$; khi ghi phải kèm đơn vị và làm tròn phù hợp độ chia nhỏ nhất.
}
\end{bt}

% ===== Câu 53 =====
% ID: AIQ_L10C3_FULL_1511
% Mức: VD
\begin{bt}
% ID: AIQ_L10C3_FULL_1511
% Mức: VD
Ba lần đo lực được $2{,}5N$, $2{,}6N$ và $2{,}7$ N. Hãy tính trung bình và nêu cách ghi kết quả.
\loigiai{
Tính trung bình cộng của ba số đo theo biểu thức $\overline F=(F_1+F_2+F_3)/3$. Kết quả là 2,6 $N$; khi ghi phải kèm đơn vị và làm tròn phù hợp độ chia nhỏ nhất.
}
\end{bt}

% ===== Câu 54 =====
% ID: AIQ_L10C3_FULL_1512
% Mức: VDC
\begin{bt}
% ID: AIQ_L10C3_FULL_1512
% Mức: VDC
Ba lần đo lực được $2{,}6N$, $2{,}7N$ và $2{,}8$ N. Hãy tính trung bình và nêu cách ghi kết quả.
\loigiai{
Tính trung bình cộng của ba số đo theo biểu thức $\overline F=(F_1+F_2+F_3)/3$. Kết quả là 2,7 $N$; khi ghi phải kèm đơn vị và làm tròn phù hợp độ chia nhỏ nhất.
}
\end{bt}

% ===== Câu 55 =====
% ID: AIQ_L10C3_FULL_1513
% Mức: NB
\dangbt{Biểu diễn và tổng hợp các vectơ lực từ số liệu}
\begin{ex}
% ID: AIQ_L10C3_FULL_1513
% Mức: NB
Hai lực vuông góc có độ lớn $F_1=3\ \mathrm{N}$ và $F_2=4\ \mathrm{N}$. Độ lớn hợp lực bằng bao nhiêu?
\choice
{\True $5\ \mathrm{N}$}
{$7\ \mathrm{N}$}
{$1\ \mathrm{N}$}
{$12\ \mathrm{N}$}
\loigiai{
Vì hai lực vuông góc nên $F=\sqrt{F_1^2+F_2^2}=5\ \mathrm{N}$.
}
\end{ex}

% ===== Câu 56 =====
% ID: AIQ_L10C3_FULL_1514
% Mức: NB
\begin{ex}
% ID: AIQ_L10C3_FULL_1514
% Mức: NB
Hai lực vuông góc có độ lớn $F_1=4\ \mathrm{N}$ và $F_2=5\ \mathrm{N}$. Độ lớn hợp lực bằng bao nhiêu?
\choice
{$9\ \mathrm{N}$}
{\True $6,4\ \mathrm{N}$}
{$1\ \mathrm{N}$}
{$20\ \mathrm{N}$}
\loigiai{
Vì hai lực vuông góc nên $F=\sqrt{F_1^2+F_2^2}=6,4\ \mathrm{N}$.
}
\end{ex}

% ===== Câu 57 =====
% ID: AIQ_L10C3_FULL_1515
% Mức: NB
\begin{ex}
% ID: AIQ_L10C3_FULL_1515
% Mức: NB
Hai lực vuông góc có độ lớn $F_1=5\ \mathrm{N}$ và $F_2=4\ \mathrm{N}$. Độ lớn hợp lực bằng bao nhiêu?
\choice
{$9\ \mathrm{N}$}
{$1\ \mathrm{N}$}
{\True $6,4\ \mathrm{N}$}
{$20\ \mathrm{N}$}
\loigiai{
Vì hai lực vuông góc nên $F=\sqrt{F_1^2+F_2^2}=6,4\ \mathrm{N}$.
}
\end{ex}

% ===== Câu 58 =====
% ID: AIQ_L10C3_FULL_1516
% Mức: TH
\begin{ex}
% ID: AIQ_L10C3_FULL_1516
% Mức: TH
Hai lực vuông góc có độ lớn $F_1=3\ \mathrm{N}$ và $F_2=5\ \mathrm{N}$. Độ lớn hợp lực bằng bao nhiêu?
\choice
{$8\ \mathrm{N}$}
{$2\ \mathrm{N}$}
{$15\ \mathrm{N}$}
{\True $5,83\ \mathrm{N}$}
\loigiai{
Vì hai lực vuông góc nên $F=\sqrt{F_1^2+F_2^2}=5,83\ \mathrm{N}$.
}
\end{ex}

% ===== Câu 59 =====
% ID: AIQ_L10C3_FULL_1517
% Mức: TH
\begin{ex}
% ID: AIQ_L10C3_FULL_1517
% Mức: TH
Hai lực vuông góc có độ lớn $F_1=4\ \mathrm{N}$ và $F_2=4\ \mathrm{N}$. Độ lớn hợp lực bằng bao nhiêu?
\choice
{\True $5,66\ \mathrm{N}$}
{$8\ \mathrm{N}$}
{$0\ \mathrm{N}$}
{$16\ \mathrm{N}$}
\loigiai{
Vì hai lực vuông góc nên $F=\sqrt{F_1^2+F_2^2}=5,66\ \mathrm{N}$.
}
\end{ex}

% ===== Câu 60 =====
% ID: AIQ_L10C3_FULL_1518
% Mức: TH
\begin{ex}
% ID: AIQ_L10C3_FULL_1518
% Mức: TH
Hai lực vuông góc có độ lớn $F_1=5\ \mathrm{N}$ và $F_2=5\ \mathrm{N}$. Độ lớn hợp lực bằng bao nhiêu?
\choice
{$10\ \mathrm{N}$}
{\True $7,07\ \mathrm{N}$}
{$0\ \mathrm{N}$}
{$25\ \mathrm{N}$}
\loigiai{
Vì hai lực vuông góc nên $F=\sqrt{F_1^2+F_2^2}=7,07\ \mathrm{N}$.
}
\end{ex}

% ===== Câu 61 =====
% ID: AIQ_L10C3_FULL_1519
% Mức: VD
\begin{ex}
% ID: AIQ_L10C3_FULL_1519
% Mức: VD
Hai lực vuông góc có độ lớn $F_1=3\ \mathrm{N}$ và $F_2=4\ \mathrm{N}$. Độ lớn hợp lực bằng bao nhiêu?
\choice
{$7\ \mathrm{N}$}
{$1\ \mathrm{N}$}
{\True $5\ \mathrm{N}$}
{$12\ \mathrm{N}$}
\loigiai{
Vì hai lực vuông góc nên $F=\sqrt{F_1^2+F_2^2}=5\ \mathrm{N}$.
}
\end{ex}

% ===== Câu 62 =====
% ID: AIQ_L10C3_FULL_1520
% Mức: VD
\begin{ex}
% ID: AIQ_L10C3_FULL_1520
% Mức: VD
Hai lực vuông góc có độ lớn $F_1=4\ \mathrm{N}$ và $F_2=5\ \mathrm{N}$. Độ lớn hợp lực bằng bao nhiêu?
\choice
{$9\ \mathrm{N}$}
{$1\ \mathrm{N}$}
{$20\ \mathrm{N}$}
{\True $6,4\ \mathrm{N}$}
\loigiai{
Vì hai lực vuông góc nên $F=\sqrt{F_1^2+F_2^2}=6,4\ \mathrm{N}$.
}
\end{ex}

% ===== Câu 63 =====
% ID: AIQ_L10C3_FULL_1521
% Mức: VD
\begin{ex}
% ID: AIQ_L10C3_FULL_1521
% Mức: VD
Hai lực vuông góc có độ lớn $F_1=5\ \mathrm{N}$ và $F_2=4\ \mathrm{N}$. Độ lớn hợp lực bằng bao nhiêu?
\choice
{\True $6,4\ \mathrm{N}$}
{$9\ \mathrm{N}$}
{$1\ \mathrm{N}$}
{$20\ \mathrm{N}$}
\loigiai{
Vì hai lực vuông góc nên $F=\sqrt{F_1^2+F_2^2}=6,4\ \mathrm{N}$.
}
\end{ex}

% ===== Câu 64 =====
% ID: AIQ_L10C3_FULL_1522
% Mức: VDC
\begin{ex}
% ID: AIQ_L10C3_FULL_1522
% Mức: VDC
Hai lực vuông góc có độ lớn $F_1=3\ \mathrm{N}$ và $F_2=5\ \mathrm{N}$. Độ lớn hợp lực bằng bao nhiêu?
\choice
{$8\ \mathrm{N}$}
{\True $5,83\ \mathrm{N}$}
{$2\ \mathrm{N}$}
{$15\ \mathrm{N}$}
\loigiai{
Vì hai lực vuông góc nên $F=\sqrt{F_1^2+F_2^2}=5,83\ \mathrm{N}$.
}
\end{ex}

% ===== Câu 65 =====
% ID: AIQ_L10C3_FULL_1523
% Mức: NB
\begin{ex}
% ID: AIQ_L10C3_FULL_1523
% Mức: NB
Hai lực vuông góc có độ lớn $3\ \mathrm{N}$ và $4\ \mathrm{N}$ trong lượt 1.
\choiceTF
{\True Hợp lực có độ lớn $5\ \mathrm{N}$.}
{\True Vectơ hợp lực là đường chéo hình bình hành.}
{Độ lớn hợp lực bằng $7\ \mathrm{N}$.}
{\True Khi vẽ tỉ xích $1\ \mathrm{cm}$ ứng với $1\ \mathrm{N}$, vectơ hợp lực dài $5\ \mathrm{cm}$.}
\loigiai{
\begin{itemchoice}
\item (Đúng) Hợp lực có độ lớn $5\ \mathrm{N}$. (Vì): vì $F=\sqrt{3^2+4^2}=5\ \mathrm{N}$. b) Đúng: hợp lực là đường chéo hình bình hành. c) Sai vì $3+4$ chỉ áp dụng khi hai lực cùng chiều. d) Đúng theo tỉ xích đã cho
\item (Đúng) Vectơ hợp lực là đường chéo hình bình hành.
\item (Sai) Độ lớn hợp lực bằng $7\ \mathrm{N}$.
\item (Đúng) Khi vẽ tỉ xích $1\ \mathrm{cm}$ ứng với $1\ \mathrm{N}$, vectơ hợp lực dài $5\ \mathrm{cm}$.
\end{itemchoice}
}
\end{ex}

% ===== Câu 66 =====
% ID: AIQ_L10C3_FULL_1524
% Mức: TH
\begin{ex}
% ID: AIQ_L10C3_FULL_1524
% Mức: TH
Hai lực vuông góc có độ lớn $3\ \mathrm{N}$ và $4\ \mathrm{N}$ trong lượt 2.
\choiceTF
{\True Hợp lực có độ lớn $5\ \mathrm{N}$.}
{\True Vectơ hợp lực là đường chéo hình bình hành.}
{Độ lớn hợp lực bằng $7\ \mathrm{N}$.}
{\True Khi vẽ tỉ xích $1\ \mathrm{cm}$ ứng với $1\ \mathrm{N}$, vectơ hợp lực dài $5\ \mathrm{cm}$.}
\loigiai{
\begin{itemchoice}
\item (Đúng) Hợp lực có độ lớn $5\ \mathrm{N}$. (Vì): vì $F=\sqrt{3^2+4^2}=5\ \mathrm{N}$. b) Đúng: hợp lực là đường chéo hình bình hành. c) Sai vì $3+4$ chỉ áp dụng khi hai lực cùng chiều. d) Đúng theo tỉ xích đã cho
\item (Đúng) Vectơ hợp lực là đường chéo hình bình hành.
\item (Sai) Độ lớn hợp lực bằng $7\ \mathrm{N}$.
\item (Đúng) Khi vẽ tỉ xích $1\ \mathrm{cm}$ ứng với $1\ \mathrm{N}$, vectơ hợp lực dài $5\ \mathrm{cm}$.
\end{itemchoice}
}
\end{ex}

% ===== Câu 67 =====
% ID: AIQ_L10C3_FULL_1525
% Mức: VD
\begin{ex}
% ID: AIQ_L10C3_FULL_1525
% Mức: VD
Hai lực vuông góc có độ lớn $3\ \mathrm{N}$ và $4\ \mathrm{N}$ trong lượt 3.
\choiceTF
{\True Hợp lực có độ lớn $5\ \mathrm{N}$.}
{\True Vectơ hợp lực là đường chéo hình bình hành.}
{Độ lớn hợp lực bằng $7\ \mathrm{N}$.}
{\True Khi vẽ tỉ xích $1\ \mathrm{cm}$ ứng với $1\ \mathrm{N}$, vectơ hợp lực dài $5\ \mathrm{cm}$.}
\loigiai{
\begin{itemchoice}
\item (Đúng) Hợp lực có độ lớn $5\ \mathrm{N}$. (Vì): vì $F=\sqrt{3^2+4^2}=5\ \mathrm{N}$. b) Đúng: hợp lực là đường chéo hình bình hành. c) Sai vì $3+4$ chỉ áp dụng khi hai lực cùng chiều. d) Đúng theo tỉ xích đã cho
\item (Đúng) Vectơ hợp lực là đường chéo hình bình hành.
\item (Sai) Độ lớn hợp lực bằng $7\ \mathrm{N}$.
\item (Đúng) Khi vẽ tỉ xích $1\ \mathrm{cm}$ ứng với $1\ \mathrm{N}$, vectơ hợp lực dài $5\ \mathrm{cm}$.
\end{itemchoice}
}
\end{ex}

% ===== Câu 68 =====
% ID: AIQ_L10C3_FULL_1526
% Mức: VD
\begin{ex}
% ID: AIQ_L10C3_FULL_1526
% Mức: VD
Hai lực vuông góc có độ lớn $3\ \mathrm{N}$ và $4\ \mathrm{N}$ trong lượt 4.
\choiceTF
{\True Hợp lực có độ lớn $5\ \mathrm{N}$.}
{\True Vectơ hợp lực là đường chéo hình bình hành.}
{Độ lớn hợp lực bằng $7\ \mathrm{N}$.}
{\True Khi vẽ tỉ xích $1\ \mathrm{cm}$ ứng với $1\ \mathrm{N}$, vectơ hợp lực dài $5\ \mathrm{cm}$.}
\loigiai{
\begin{itemchoice}
\item (Đúng) Hợp lực có độ lớn $5\ \mathrm{N}$. (Vì): vì $F=\sqrt{3^2+4^2}=5\ \mathrm{N}$. b) Đúng: hợp lực là đường chéo hình bình hành. c) Sai vì $3+4$ chỉ áp dụng khi hai lực cùng chiều. d) Đúng theo tỉ xích đã cho
\item (Đúng) Vectơ hợp lực là đường chéo hình bình hành.
\item (Sai) Độ lớn hợp lực bằng $7\ \mathrm{N}$.
\item (Đúng) Khi vẽ tỉ xích $1\ \mathrm{cm}$ ứng với $1\ \mathrm{N}$, vectơ hợp lực dài $5\ \mathrm{cm}$.
\end{itemchoice}
}
\end{ex}

% ===== Câu 69 =====
% ID: AIQ_L10C3_FULL_1527
% Mức: VDC
\begin{ex}
% ID: AIQ_L10C3_FULL_1527
% Mức: VDC
Hai lực vuông góc có độ lớn $3\ \mathrm{N}$ và $4\ \mathrm{N}$ trong lượt 5.
\choiceTF
{\True Hợp lực có độ lớn $5\ \mathrm{N}$.}
{\True Vectơ hợp lực là đường chéo hình bình hành.}
{Độ lớn hợp lực bằng $7\ \mathrm{N}$.}
{\True Khi vẽ tỉ xích $1\ \mathrm{cm}$ ứng với $1\ \mathrm{N}$, vectơ hợp lực dài $5\ \mathrm{cm}$.}
\loigiai{
\begin{itemchoice}
\item (Đúng) Hợp lực có độ lớn $5\ \mathrm{N}$. (Vì): vì $F=\sqrt{3^2+4^2}=5\ \mathrm{N}$. b) Đúng: hợp lực là đường chéo hình bình hành. c) Sai vì $3+4$ chỉ áp dụng khi hai lực cùng chiều. d) Đúng theo tỉ xích đã cho
\item (Đúng) Vectơ hợp lực là đường chéo hình bình hành.
\item (Sai) Độ lớn hợp lực bằng $7\ \mathrm{N}$.
\item (Đúng) Khi vẽ tỉ xích $1\ \mathrm{cm}$ ứng với $1\ \mathrm{N}$, vectơ hợp lực dài $5\ \mathrm{cm}$.
\end{itemchoice}
}
\end{ex}

% ===== Câu 70 =====
% ID: AIQ_L10C3_FULL_1528
% Mức: TH
\begin{ex}
% ID: AIQ_L10C3_FULL_1528
% Mức: TH
Hai lực vuông góc có độ lớn $3N$ và $4$ N. Tính độ lớn hợp lực theo $N$, làm tròn đến hai chữ số thập phân.
\shortans{5}
\loigiai{
Hai lực vuông góc nên áp dụng định lí Pythagore: $F=\sqrt{F_1^2+F_2^2}=5\ \mathrm{N}$. Đáp án không ghi đơn vị.
}
\end{ex}

% ===== Câu 71 =====
% ID: AIQ_L10C3_FULL_1529
% Mức: TH
\begin{ex}
% ID: AIQ_L10C3_FULL_1529
% Mức: TH
Hai lực vuông góc có độ lớn $4N$ và $4$ N. Tính độ lớn hợp lực theo $N$, làm tròn đến hai chữ số thập phân.
\shortans{5,66}
\loigiai{
Hai lực vuông góc nên áp dụng định lí Pythagore: $F=\sqrt{F_1^2+F_2^2}=5,66\ \mathrm{N}$. Đáp án không ghi đơn vị.
}
\end{ex}

% ===== Câu 72 =====
% ID: AIQ_L10C3_FULL_1530
% Mức: VD
\begin{ex}
% ID: AIQ_L10C3_FULL_1530
% Mức: VD
Hai lực vuông góc có độ lớn $5N$ và $4$ N. Tính độ lớn hợp lực theo $N$, làm tròn đến hai chữ số thập phân.
\shortans{6,4}
\loigiai{
Hai lực vuông góc nên áp dụng định lí Pythagore: $F=\sqrt{F_1^2+F_2^2}=6,4\ \mathrm{N}$. Đáp án không ghi đơn vị.
}
\end{ex}

% ===== Câu 73 =====
% ID: AIQ_L10C3_FULL_1531
% Mức: VD
\begin{ex}
% ID: AIQ_L10C3_FULL_1531
% Mức: VD
Hai lực vuông góc có độ lớn $6N$ và $4$ N. Tính độ lớn hợp lực theo $N$, làm tròn đến hai chữ số thập phân.
\shortans{7,21}
\loigiai{
Hai lực vuông góc nên áp dụng định lí Pythagore: $F=\sqrt{F_1^2+F_2^2}=7,21\ \mathrm{N}$. Đáp án không ghi đơn vị.
}
\end{ex}

% ===== Câu 74 =====
% ID: AIQ_L10C3_FULL_1532
% Mức: VD
\begin{ex}
% ID: AIQ_L10C3_FULL_1532
% Mức: VD
Hai lực vuông góc có độ lớn $7N$ và $4$ N. Tính độ lớn hợp lực theo $N$, làm tròn đến hai chữ số thập phân.
\shortans{8,06}
\loigiai{
Hai lực vuông góc nên áp dụng định lí Pythagore: $F=\sqrt{F_1^2+F_2^2}=8,06\ \mathrm{N}$. Đáp án không ghi đơn vị.
}
\end{ex}

% ===== Câu 75 =====
% ID: AIQ_L10C3_FULL_1533
% Mức: VDC
\begin{ex}
% ID: AIQ_L10C3_FULL_1533
% Mức: VDC
Hai lực vuông góc có độ lớn $8N$ và $4$ N. Tính độ lớn hợp lực theo $N$, làm tròn đến hai chữ số thập phân.
\shortans{8,94}
\loigiai{
Hai lực vuông góc nên áp dụng định lí Pythagore: $F=\sqrt{F_1^2+F_2^2}=8,94\ \mathrm{N}$. Đáp án không ghi đơn vị.
}
\end{ex}

% ===== Câu 76 =====
% ID: AIQ_L10C3_FULL_1534
% Mức: TH
\begin{bt}
% ID: AIQ_L10C3_FULL_1534
% Mức: TH
Hai lực vuông góc có độ lớn $5N$ và $12$ N. Hãy biểu diễn theo tỉ xích $1$ cm ứng với $1N$ và xác định hợp lực.
\loigiai{
Vẽ hai vectơ lực vuông góc theo tỉ xích đã cho. Dựng hình bình hành; đường chéo xuất phát từ điểm đồng quy biểu diễn hợp lực. Theo định lí Pythagore thu được $13\,\text{N}$.
}
\end{bt}

% ===== Câu 77 =====
% ID: AIQ_L10C3_FULL_1535
% Mức: TH
\begin{bt}
% ID: AIQ_L10C3_FULL_1535
% Mức: TH
Hai lực vuông góc có độ lớn $6N$ và $12$ N. Hãy biểu diễn theo tỉ xích $1$ cm ứng với $1N$ và xác định hợp lực.
\loigiai{
Vẽ hai vectơ lực vuông góc theo tỉ xích đã cho. Dựng hình bình hành; đường chéo xuất phát từ điểm đồng quy biểu diễn hợp lực. Theo định lí Pythagore thu được $13,42\,\text{N}$.
}
\end{bt}

% ===== Câu 78 =====
% ID: AIQ_L10C3_FULL_1536
% Mức: VD
\begin{bt}
% ID: AIQ_L10C3_FULL_1536
% Mức: VD
Hai lực vuông góc có độ lớn $7N$ và $12$ N. Hãy biểu diễn theo tỉ xích $1$ cm ứng với $1N$ và xác định hợp lực.
\loigiai{
Vẽ hai vectơ lực vuông góc theo tỉ xích đã cho. Dựng hình bình hành; đường chéo xuất phát từ điểm đồng quy biểu diễn hợp lực. Theo định lí Pythagore thu được $13,89\,\text{N}$.
}
\end{bt}

% ===== Câu 79 =====
% ID: AIQ_L10C3_FULL_1537
% Mức: VD
\begin{bt}
% ID: AIQ_L10C3_FULL_1537
% Mức: VD
Hai lực vuông góc có độ lớn $8N$ và $12$ N. Hãy biểu diễn theo tỉ xích $1$ cm ứng với $1N$ và xác định hợp lực.
\loigiai{
Vẽ hai vectơ lực vuông góc theo tỉ xích đã cho. Dựng hình bình hành; đường chéo xuất phát từ điểm đồng quy biểu diễn hợp lực. Theo định lí Pythagore thu được $14,42\,\text{N}$.
}
\end{bt}

% ===== Câu 80 =====
% ID: AIQ_L10C3_FULL_1538
% Mức: VD
\begin{bt}
% ID: AIQ_L10C3_FULL_1538
% Mức: VD
Hai lực vuông góc có độ lớn $9N$ và $12$ N. Hãy biểu diễn theo tỉ xích $1$ cm ứng với $1N$ và xác định hợp lực.
\loigiai{
Vẽ hai vectơ lực vuông góc theo tỉ xích đã cho. Dựng hình bình hành; đường chéo xuất phát từ điểm đồng quy biểu diễn hợp lực. Theo định lí Pythagore thu được $15\,\text{N}$.
}
\end{bt}

% ===== Câu 81 =====
% ID: AIQ_L10C3_FULL_1539
% Mức: VDC
\begin{bt}
% ID: AIQ_L10C3_FULL_1539
% Mức: VDC
Hai lực vuông góc có độ lớn $10N$ và $12$ N. Hãy biểu diễn theo tỉ xích $1$ cm ứng với $1N$ và xác định hợp lực.
\loigiai{
Vẽ hai vectơ lực vuông góc theo tỉ xích đã cho. Dựng hình bình hành; đường chéo xuất phát từ điểm đồng quy biểu diễn hợp lực. Theo định lí Pythagore thu được $15,62\,\text{N}$.
}
\end{bt}

% ===== Câu 82 =====
% ID: AIQ_L10C3_FULL_1540
% Mức: NB
\dangbt{Kiểm chứng quy tắc hình bình hành}
\begin{ex}
% ID: AIQ_L10C3_FULL_1540
% Mức: NB
Hai lực đồng quy cùng độ lớn $2\ \mathrm{N}$ hợp nhau góc $0^\circ$. Theo quy tắc hình bình hành, hợp lực có độ lớn bao nhiêu?
\choice
{\True $4\ \mathrm{N}$}
{$1\ \mathrm{N}$}
{$2,5\ \mathrm{N}$}
{$5\ \mathrm{N}$}
\loigiai{
$F=\sqrt{F^2+F^2+2F^2\cos 0^\circ}=2F\cos(0^\circ/2)=4\ \mathrm{N}$.
}
\end{ex}

% ===== Câu 83 =====
% ID: AIQ_L10C3_FULL_1541
% Mức: NB
\begin{ex}
% ID: AIQ_L10C3_FULL_1541
% Mức: NB
Hai lực đồng quy cùng độ lớn $3\ \mathrm{N}$ hợp nhau góc $60^\circ$. Theo quy tắc hình bình hành, hợp lực có độ lớn bao nhiêu?
\choice
{$1,5\ \mathrm{N}$}
{\True $5,2\ \mathrm{N}$}
{$3,75\ \mathrm{N}$}
{$7,5\ \mathrm{N}$}
\loigiai{
$F=\sqrt{F^2+F^2+2F^2\cos 60^\circ}=2F\cos(60^\circ/2)=5,2\ \mathrm{N}$.
}
\end{ex}

% ===== Câu 84 =====
% ID: AIQ_L10C3_FULL_1542
% Mức: NB
\begin{ex}
% ID: AIQ_L10C3_FULL_1542
% Mức: NB
Hai lực đồng quy cùng độ lớn $4\ \mathrm{N}$ hợp nhau góc $90^\circ$. Theo quy tắc hình bình hành, hợp lực có độ lớn bao nhiêu?
\choice
{$2\ \mathrm{N}$}
{$5\ \mathrm{N}$}
{\True $5,66\ \mathrm{N}$}
{$10\ \mathrm{N}$}
\loigiai{
$F=\sqrt{F^2+F^2+2F^2\cos 90^\circ}=2F\cos(90^\circ/2)=5,66\ \mathrm{N}$.
}
\end{ex}

% ===== Câu 85 =====
% ID: AIQ_L10C3_FULL_1543
% Mức: TH
\begin{ex}
% ID: AIQ_L10C3_FULL_1543
% Mức: TH
Hai lực đồng quy cùng độ lớn $5\ \mathrm{N}$ hợp nhau góc $120^\circ$. Theo quy tắc hình bình hành, hợp lực có độ lớn bao nhiêu?
\choice
{$2,5\ \mathrm{N}$}
{$6,25\ \mathrm{N}$}
{$12,5\ \mathrm{N}$}
{\True $5\ \mathrm{N}$}
\loigiai{
$F=\sqrt{F^2+F^2+2F^2\cos 120^\circ}=2F\cos(120^\circ/2)=5\ \mathrm{N}$.
}
\end{ex}

% ===== Câu 86 =====
% ID: AIQ_L10C3_FULL_1544
% Mức: TH
\begin{ex}
% ID: AIQ_L10C3_FULL_1544
% Mức: TH
Hai lực đồng quy cùng độ lớn $2\ \mathrm{N}$ hợp nhau góc $180^\circ$. Theo quy tắc hình bình hành, hợp lực có độ lớn bao nhiêu?
\choice
{\True $0\ \mathrm{N}$}
{$1\ \mathrm{N}$}
{$2,5\ \mathrm{N}$}
{$5\ \mathrm{N}$}
\loigiai{
$F=\sqrt{F^2+F^2+2F^2\cos 180^\circ}=2F\cos(180^\circ/2)=0\ \mathrm{N}$.
}
\end{ex}

% ===== Câu 87 =====
% ID: AIQ_L10C3_FULL_1545
% Mức: TH
\begin{ex}
% ID: AIQ_L10C3_FULL_1545
% Mức: TH
Hai lực đồng quy cùng độ lớn $3\ \mathrm{N}$ hợp nhau góc $0^\circ$. Theo quy tắc hình bình hành, hợp lực có độ lớn bao nhiêu?
\choice
{$1,5\ \mathrm{N}$}
{\True $6\ \mathrm{N}$}
{$3,75\ \mathrm{N}$}
{$7,5\ \mathrm{N}$}
\loigiai{
$F=\sqrt{F^2+F^2+2F^2\cos 0^\circ}=2F\cos(0^\circ/2)=6\ \mathrm{N}$.
}
\end{ex}

% ===== Câu 88 =====
% ID: AIQ_L10C3_FULL_1546
% Mức: VD
\begin{ex}
% ID: AIQ_L10C3_FULL_1546
% Mức: VD
Hai lực đồng quy cùng độ lớn $4\ \mathrm{N}$ hợp nhau góc $60^\circ$. Theo quy tắc hình bình hành, hợp lực có độ lớn bao nhiêu?
\choice
{$2\ \mathrm{N}$}
{$5\ \mathrm{N}$}
{\True $6,93\ \mathrm{N}$}
{$10\ \mathrm{N}$}
\loigiai{
$F=\sqrt{F^2+F^2+2F^2\cos 60^\circ}=2F\cos(60^\circ/2)=6,93\ \mathrm{N}$.
}
\end{ex}

% ===== Câu 89 =====
% ID: AIQ_L10C3_FULL_1547
% Mức: VD
\begin{ex}
% ID: AIQ_L10C3_FULL_1547
% Mức: VD
Hai lực đồng quy cùng độ lớn $5\ \mathrm{N}$ hợp nhau góc $90^\circ$. Theo quy tắc hình bình hành, hợp lực có độ lớn bao nhiêu?
\choice
{$2,5\ \mathrm{N}$}
{$6,25\ \mathrm{N}$}
{$12,5\ \mathrm{N}$}
{\True $7,07\ \mathrm{N}$}
\loigiai{
$F=\sqrt{F^2+F^2+2F^2\cos 90^\circ}=2F\cos(90^\circ/2)=7,07\ \mathrm{N}$.
}
\end{ex}

% ===== Câu 90 =====
% ID: AIQ_L10C3_FULL_1548
% Mức: VD
\begin{ex}
% ID: AIQ_L10C3_FULL_1548
% Mức: VD
Hai lực đồng quy cùng độ lớn $2\ \mathrm{N}$ hợp nhau góc $120^\circ$. Theo quy tắc hình bình hành, hợp lực có độ lớn bao nhiêu?
\choice
{\True $2\ \mathrm{N}$}
{$1\ \mathrm{N}$}
{$2,5\ \mathrm{N}$}
{$5\ \mathrm{N}$}
\loigiai{
$F=\sqrt{F^2+F^2+2F^2\cos 120^\circ}=2F\cos(120^\circ/2)=2\ \mathrm{N}$.
}
\end{ex}

% ===== Câu 91 =====
% ID: AIQ_L10C3_FULL_1549
% Mức: VDC
\begin{ex}
% ID: AIQ_L10C3_FULL_1549
% Mức: VDC
Hai lực đồng quy cùng độ lớn $3\ \mathrm{N}$ hợp nhau góc $180^\circ$. Theo quy tắc hình bình hành, hợp lực có độ lớn bao nhiêu?
\choice
{$1,5\ \mathrm{N}$}
{\True $0\ \mathrm{N}$}
{$3,75\ \mathrm{N}$}
{$7,5\ \mathrm{N}$}
\loigiai{
$F=\sqrt{F^2+F^2+2F^2\cos 180^\circ}=2F\cos(180^\circ/2)=0\ \mathrm{N}$.
}
\end{ex}

% ===== Câu 92 =====
% ID: AIQ_L10C3_FULL_1550
% Mức: NB
\begin{ex}
% ID: AIQ_L10C3_FULL_1550
% Mức: NB
Kiểm chứng quy tắc hình bình hành ở lượt 1.
\choiceTF
{\True Hai cạnh kề biểu diễn hai lực thành phần.}
{\True Đường chéo xuất phát từ điểm đồng quy biểu diễn hợp lực.}
{Hai lực cùng chiều thì hợp lực có độ lớn bằng hiệu hai lực.}
{\True Hai lực ngược chiều thì hợp lực có độ lớn bằng trị tuyệt đối hiệu hai lực.}
\loigiai{
\begin{itemchoice}
\item (Đúng) Hai cạnh kề biểu diễn hai lực thành phần. (Vì): hai cạnh kề biểu diễn hai lực thành phần. b) Đúng: đường chéo từ điểm đồng quy biểu diễn hợp lực. c) Sai: cùng chiều thì hợp lực bằng tổng độ lớn. d) Đúng: ngược chiều thì hợp lực bằng trị tuyệt đối hiệu
\item (Đúng) Đường chéo xuất phát từ điểm đồng quy biểu diễn hợp lực.
\item (Sai) Hai lực cùng chiều thì hợp lực có độ lớn bằng hiệu hai lực.
\item (Đúng) Hai lực ngược chiều thì hợp lực có độ lớn bằng trị tuyệt đối hiệu hai lực.
\end{itemchoice}
}
\end{ex}

% ===== Câu 93 =====
% ID: AIQ_L10C3_FULL_1551
% Mức: TH
\begin{ex}
% ID: AIQ_L10C3_FULL_1551
% Mức: TH
Kiểm chứng quy tắc hình bình hành ở lượt 2.
\choiceTF
{\True Hai cạnh kề biểu diễn hai lực thành phần.}
{\True Đường chéo xuất phát từ điểm đồng quy biểu diễn hợp lực.}
{Hai lực cùng chiều thì hợp lực có độ lớn bằng hiệu hai lực.}
{\True Hai lực ngược chiều thì hợp lực có độ lớn bằng trị tuyệt đối hiệu hai lực.}
\loigiai{
\begin{itemchoice}
\item (Đúng) Hai cạnh kề biểu diễn hai lực thành phần. (Vì): hai cạnh kề biểu diễn hai lực thành phần. b) Đúng: đường chéo từ điểm đồng quy biểu diễn hợp lực. c) Sai: cùng chiều thì hợp lực bằng tổng độ lớn. d) Đúng: ngược chiều thì hợp lực bằng trị tuyệt đối hiệu
\item (Đúng) Đường chéo xuất phát từ điểm đồng quy biểu diễn hợp lực.
\item (Sai) Hai lực cùng chiều thì hợp lực có độ lớn bằng hiệu hai lực.
\item (Đúng) Hai lực ngược chiều thì hợp lực có độ lớn bằng trị tuyệt đối hiệu hai lực.
\end{itemchoice}
}
\end{ex}

% ===== Câu 94 =====
% ID: AIQ_L10C3_FULL_1552
% Mức: VD
\begin{ex}
% ID: AIQ_L10C3_FULL_1552
% Mức: VD
Kiểm chứng quy tắc hình bình hành ở lượt 3.
\choiceTF
{\True Hai cạnh kề biểu diễn hai lực thành phần.}
{\True Đường chéo xuất phát từ điểm đồng quy biểu diễn hợp lực.}
{Hai lực cùng chiều thì hợp lực có độ lớn bằng hiệu hai lực.}
{\True Hai lực ngược chiều thì hợp lực có độ lớn bằng trị tuyệt đối hiệu hai lực.}
\loigiai{
\begin{itemchoice}
\item (Đúng) Hai cạnh kề biểu diễn hai lực thành phần. (Vì): hai cạnh kề biểu diễn hai lực thành phần. b) Đúng: đường chéo từ điểm đồng quy biểu diễn hợp lực. c) Sai: cùng chiều thì hợp lực bằng tổng độ lớn. d) Đúng: ngược chiều thì hợp lực bằng trị tuyệt đối hiệu
\item (Đúng) Đường chéo xuất phát từ điểm đồng quy biểu diễn hợp lực.
\item (Sai) Hai lực cùng chiều thì hợp lực có độ lớn bằng hiệu hai lực.
\item (Đúng) Hai lực ngược chiều thì hợp lực có độ lớn bằng trị tuyệt đối hiệu hai lực.
\end{itemchoice}
}
\end{ex}

% ===== Câu 95 =====
% ID: AIQ_L10C3_FULL_1553
% Mức: VD
\begin{ex}
% ID: AIQ_L10C3_FULL_1553
% Mức: VD
Kiểm chứng quy tắc hình bình hành ở lượt 4.
\choiceTF
{\True Hai cạnh kề biểu diễn hai lực thành phần.}
{\True Đường chéo xuất phát từ điểm đồng quy biểu diễn hợp lực.}
{Hai lực cùng chiều thì hợp lực có độ lớn bằng hiệu hai lực.}
{\True Hai lực ngược chiều thì hợp lực có độ lớn bằng trị tuyệt đối hiệu hai lực.}
\loigiai{
\begin{itemchoice}
\item (Đúng) Hai cạnh kề biểu diễn hai lực thành phần. (Vì): hai cạnh kề biểu diễn hai lực thành phần. b) Đúng: đường chéo từ điểm đồng quy biểu diễn hợp lực. c) Sai: cùng chiều thì hợp lực bằng tổng độ lớn. d) Đúng: ngược chiều thì hợp lực bằng trị tuyệt đối hiệu
\item (Đúng) Đường chéo xuất phát từ điểm đồng quy biểu diễn hợp lực.
\item (Sai) Hai lực cùng chiều thì hợp lực có độ lớn bằng hiệu hai lực.
\item (Đúng) Hai lực ngược chiều thì hợp lực có độ lớn bằng trị tuyệt đối hiệu hai lực.
\end{itemchoice}
}
\end{ex}

% ===== Câu 96 =====
% ID: AIQ_L10C3_FULL_1554
% Mức: VDC
\begin{ex}
% ID: AIQ_L10C3_FULL_1554
% Mức: VDC
Kiểm chứng quy tắc hình bình hành ở lượt 5.
\choiceTF
{\True Hai cạnh kề biểu diễn hai lực thành phần.}
{\True Đường chéo xuất phát từ điểm đồng quy biểu diễn hợp lực.}
{Hai lực cùng chiều thì hợp lực có độ lớn bằng hiệu hai lực.}
{\True Hai lực ngược chiều thì hợp lực có độ lớn bằng trị tuyệt đối hiệu hai lực.}
\loigiai{
\begin{itemchoice}
\item (Đúng) Hai cạnh kề biểu diễn hai lực thành phần. (Vì): hai cạnh kề biểu diễn hai lực thành phần. b) Đúng: đường chéo từ điểm đồng quy biểu diễn hợp lực. c) Sai: cùng chiều thì hợp lực bằng tổng độ lớn. d) Đúng: ngược chiều thì hợp lực bằng trị tuyệt đối hiệu
\item (Đúng) Đường chéo xuất phát từ điểm đồng quy biểu diễn hợp lực.
\item (Sai) Hai lực cùng chiều thì hợp lực có độ lớn bằng hiệu hai lực.
\item (Đúng) Hai lực ngược chiều thì hợp lực có độ lớn bằng trị tuyệt đối hiệu hai lực.
\end{itemchoice}
}
\end{ex}

% ===== Câu 97 =====
% ID: AIQ_L10C3_FULL_1555
% Mức: TH
\begin{ex}
% ID: AIQ_L10C3_FULL_1555
% Mức: TH
Hai lực bằng nhau, mỗi lực $2N$, hợp nhau góc $60^\circ$. Tính hợp lực theo $N$, làm tròn đến hai chữ số thập phân.
\shortans{3,46}
\loigiai{
Với hai lực bằng nhau hợp góc $60^\circ$: $R=2F\cos30^\circ=F\sqrt3=3,46\ \mathrm{N}$. Đáp án không ghi đơn vị.
}
\end{ex}

% ===== Câu 98 =====
% ID: AIQ_L10C3_FULL_1556
% Mức: TH
\begin{ex}
% ID: AIQ_L10C3_FULL_1556
% Mức: TH
Hai lực bằng nhau, mỗi lực $3N$, hợp nhau góc $60^\circ$. Tính hợp lực theo $N$, làm tròn đến hai chữ số thập phân.
\shortans{5,2}
\loigiai{
Với hai lực bằng nhau hợp góc $60^\circ$: $R=2F\cos30^\circ=F\sqrt3=5,2\ \mathrm{N}$. Đáp án không ghi đơn vị.
}
\end{ex}

% ===== Câu 99 =====
% ID: AIQ_L10C3_FULL_1557
% Mức: VD
\begin{ex}
% ID: AIQ_L10C3_FULL_1557
% Mức: VD
Hai lực bằng nhau, mỗi lực $4N$, hợp nhau góc $60^\circ$. Tính hợp lực theo $N$, làm tròn đến hai chữ số thập phân.
\shortans{6,93}
\loigiai{
Với hai lực bằng nhau hợp góc $60^\circ$: $R=2F\cos30^\circ=F\sqrt3=6,93\ \mathrm{N}$. Đáp án không ghi đơn vị.
}
\end{ex}

% ===== Câu 100 =====
% ID: AIQ_L10C3_FULL_1558
% Mức: VD
\begin{ex}
% ID: AIQ_L10C3_FULL_1558
% Mức: VD
Hai lực bằng nhau, mỗi lực $5N$, hợp nhau góc $60^\circ$. Tính hợp lực theo $N$, làm tròn đến hai chữ số thập phân.
\shortans{8,66}
\loigiai{
Với hai lực bằng nhau hợp góc $60^\circ$: $R=2F\cos30^\circ=F\sqrt3=8,66\ \mathrm{N}$. Đáp án không ghi đơn vị.
}
\end{ex}

% ===== Câu 101 =====
% ID: AIQ_L10C3_FULL_1559
% Mức: VD
\begin{ex}
% ID: AIQ_L10C3_FULL_1559
% Mức: VD
Hai lực bằng nhau, mỗi lực $6N$, hợp nhau góc $60^\circ$. Tính hợp lực theo $N$, làm tròn đến hai chữ số thập phân.
\shortans{10,39}
\loigiai{
Với hai lực bằng nhau hợp góc $60^\circ$: $R=2F\cos30^\circ=F\sqrt3=10,39\ \mathrm{N}$. Đáp án không ghi đơn vị.
}
\end{ex}

% ===== Câu 102 =====
% ID: AIQ_L10C3_FULL_1560
% Mức: VDC
\begin{ex}
% ID: AIQ_L10C3_FULL_1560
% Mức: VDC
Hai lực bằng nhau, mỗi lực $7N$, hợp nhau góc $60^\circ$. Tính hợp lực theo $N$, làm tròn đến hai chữ số thập phân.
\shortans{12,12}
\loigiai{
Với hai lực bằng nhau hợp góc $60^\circ$: $R=2F\cos30^\circ=F\sqrt3=12,12\ \mathrm{N}$. Đáp án không ghi đơn vị.
}
\end{ex}

% ===== Câu 103 =====
% ID: AIQ_L10C3_FULL_1561
% Mức: TH
\begin{bt}
% ID: AIQ_L10C3_FULL_1561
% Mức: TH
Hai lực bằng nhau, mỗi lực $4N$, hợp nhau góc $120^\circ$. Dùng quy tắc hình bình hành tính hợp lực và nêu phương của nó.
\loigiai{
Với hai lực bằng nhau hợp góc $120^\circ$, $R=2F\cos60^\circ=F$, nên $4\,\text{N}$. Hợp lực nằm trên phân giác trong của góc giữa hai lực.
}
\end{bt}

% ===== Câu 104 =====
% ID: AIQ_L10C3_FULL_1562
% Mức: TH
\begin{bt}
% ID: AIQ_L10C3_FULL_1562
% Mức: TH
Hai lực bằng nhau, mỗi lực $5N$, hợp nhau góc $120^\circ$. Dùng quy tắc hình bình hành tính hợp lực và nêu phương của nó.
\loigiai{
Với hai lực bằng nhau hợp góc $120^\circ$, $R=2F\cos60^\circ=F$, nên $5\,\text{N}$. Hợp lực nằm trên phân giác trong của góc giữa hai lực.
}
\end{bt}

% ===== Câu 105 =====
% ID: AIQ_L10C3_FULL_1563
% Mức: VD
\begin{bt}
% ID: AIQ_L10C3_FULL_1563
% Mức: VD
Hai lực bằng nhau, mỗi lực $6N$, hợp nhau góc $120^\circ$. Dùng quy tắc hình bình hành tính hợp lực và nêu phương của nó.
\loigiai{
Với hai lực bằng nhau hợp góc $120^\circ$, $R=2F\cos60^\circ=F$, nên $6\,\text{N}$. Hợp lực nằm trên phân giác trong của góc giữa hai lực.
}
\end{bt}

% ===== Câu 106 =====
% ID: AIQ_L10C3_FULL_1564
% Mức: VD
\begin{bt}
% ID: AIQ_L10C3_FULL_1564
% Mức: VD
Hai lực bằng nhau, mỗi lực $7N$, hợp nhau góc $120^\circ$. Dùng quy tắc hình bình hành tính hợp lực và nêu phương của nó.
\loigiai{
Với hai lực bằng nhau hợp góc $120^\circ$, $R=2F\cos60^\circ=F$, nên $7\,\text{N}$. Hợp lực nằm trên phân giác trong của góc giữa hai lực.
}
\end{bt}

% ===== Câu 107 =====
% ID: AIQ_L10C3_FULL_1565
% Mức: VD
\begin{bt}
% ID: AIQ_L10C3_FULL_1565
% Mức: VD
Hai lực bằng nhau, mỗi lực $8N$, hợp nhau góc $120^\circ$. Dùng quy tắc hình bình hành tính hợp lực và nêu phương của nó.
\loigiai{
Với hai lực bằng nhau hợp góc $120^\circ$, $R=2F\cos60^\circ=F$, nên $8\,\text{N}$. Hợp lực nằm trên phân giác trong của góc giữa hai lực.
}
\end{bt}

% ===== Câu 108 =====
% ID: AIQ_L10C3_FULL_1566
% Mức: VDC
\begin{bt}
% ID: AIQ_L10C3_FULL_1566
% Mức: VDC
Hai lực bằng nhau, mỗi lực $9N$, hợp nhau góc $120^\circ$. Dùng quy tắc hình bình hành tính hợp lực và nêu phương của nó.
\loigiai{
Với hai lực bằng nhau hợp góc $120^\circ$, $R=2F\cos60^\circ=F$, nên $9\,\text{N}$. Hợp lực nằm trên phân giác trong của góc giữa hai lực.
}
\end{bt}

% ===== Câu 109 =====
% ID: AIQ_L10C3_FULL_1567
% Mức: NB
\dangbt{Đánh giá sai số và nguyên nhân sai lệch của thí nghiệm}
\begin{ex}
% ID: AIQ_L10C3_FULL_1567
% Mức: NB
Giá trị thực nghiệm là $4\ \mathrm{N}$, giá trị lí thuyết là $4,2\ \mathrm{N}$. Sai số tương đối gần nhất bằng bao nhiêu?
\choice
{\True $4,76\%$}
{$0,2\%$}
{$9,76\%$}
{$2,38\%$}
\loigiai{
$\delta=|F_{tn}-F_{lt}|/F_{lt}\cdot100\%=4,76\%$.
}
\end{ex}

% ===== Câu 110 =====
% ID: AIQ_L10C3_FULL_1568
% Mức: NB
\begin{ex}
% ID: AIQ_L10C3_FULL_1568
% Mức: NB
Giá trị thực nghiệm là $4,1\ \mathrm{N}$, giá trị lí thuyết là $4,3\ \mathrm{N}$. Sai số tương đối gần nhất bằng bao nhiêu?
\choice
{$0,2\%$}
{\True $4,65\%$}
{$9,65\%$}
{$2,33\%$}
\loigiai{
$\delta=|F_{tn}-F_{lt}|/F_{lt}\cdot100\%=4,65\%$.
}
\end{ex}

% ===== Câu 111 =====
% ID: AIQ_L10C3_FULL_1569
% Mức: NB
\begin{ex}
% ID: AIQ_L10C3_FULL_1569
% Mức: NB
Giá trị thực nghiệm là $4,2\ \mathrm{N}$, giá trị lí thuyết là $4,4\ \mathrm{N}$. Sai số tương đối gần nhất bằng bao nhiêu?
\choice
{$0,2\%$}
{$9,55\%$}
{\True $4,55\%$}
{$2,27\%$}
\loigiai{
$\delta=|F_{tn}-F_{lt}|/F_{lt}\cdot100\%=4,55\%$.
}
\end{ex}

% ===== Câu 112 =====
% ID: AIQ_L10C3_FULL_1570
% Mức: TH
\begin{ex}
% ID: AIQ_L10C3_FULL_1570
% Mức: TH
Giá trị thực nghiệm là $4,3\ \mathrm{N}$, giá trị lí thuyết là $4,5\ \mathrm{N}$. Sai số tương đối gần nhất bằng bao nhiêu?
\choice
{$0,2\%$}
{$9,44\%$}
{$2,22\%$}
{\True $4,44\%$}
\loigiai{
$\delta=|F_{tn}-F_{lt}|/F_{lt}\cdot100\%=4,44\%$.
}
\end{ex}

% ===== Câu 113 =====
% ID: AIQ_L10C3_FULL_1571
% Mức: TH
\begin{ex}
% ID: AIQ_L10C3_FULL_1571
% Mức: TH
Giá trị thực nghiệm là $4,4\ \mathrm{N}$, giá trị lí thuyết là $4,6\ \mathrm{N}$. Sai số tương đối gần nhất bằng bao nhiêu?
\choice
{\True $4,35\%$}
{$0,2\%$}
{$9,35\%$}
{$2,17\%$}
\loigiai{
$\delta=|F_{tn}-F_{lt}|/F_{lt}\cdot100\%=4,35\%$.
}
\end{ex}

% ===== Câu 114 =====
% ID: AIQ_L10C3_FULL_1572
% Mức: TH
\begin{ex}
% ID: AIQ_L10C3_FULL_1572
% Mức: TH
Giá trị thực nghiệm là $4,5\ \mathrm{N}$, giá trị lí thuyết là $4,7\ \mathrm{N}$. Sai số tương đối gần nhất bằng bao nhiêu?
\choice
{$0,2\%$}
{\True $4,26\%$}
{$9,26\%$}
{$2,13\%$}
\loigiai{
$\delta=|F_{tn}-F_{lt}|/F_{lt}\cdot100\%=4,26\%$.
}
\end{ex}

% ===== Câu 115 =====
% ID: AIQ_L10C3_FULL_1573
% Mức: VD
\begin{ex}
% ID: AIQ_L10C3_FULL_1573
% Mức: VD
Giá trị thực nghiệm là $4,6\ \mathrm{N}$, giá trị lí thuyết là $4,8\ \mathrm{N}$. Sai số tương đối gần nhất bằng bao nhiêu?
\choice
{$0,2\%$}
{$9,17\%$}
{\True $4,17\%$}
{$2,08\%$}
\loigiai{
$\delta=|F_{tn}-F_{lt}|/F_{lt}\cdot100\%=4,17\%$.
}
\end{ex}

% ===== Câu 116 =====
% ID: AIQ_L10C3_FULL_1574
% Mức: VD
\begin{ex}
% ID: AIQ_L10C3_FULL_1574
% Mức: VD
Giá trị thực nghiệm là $4,7\ \mathrm{N}$, giá trị lí thuyết là $4,9\ \mathrm{N}$. Sai số tương đối gần nhất bằng bao nhiêu?
\choice
{$0,2\%$}
{$9,08\%$}
{$2,04\%$}
{\True $4,08\%$}
\loigiai{
$\delta=|F_{tn}-F_{lt}|/F_{lt}\cdot100\%=4,08\%$.
}
\end{ex}

% ===== Câu 117 =====
% ID: AIQ_L10C3_FULL_1575
% Mức: VD
\begin{ex}
% ID: AIQ_L10C3_FULL_1575
% Mức: VD
Giá trị thực nghiệm là $4,8\ \mathrm{N}$, giá trị lí thuyết là $5\ \mathrm{N}$. Sai số tương đối gần nhất bằng bao nhiêu?
\choice
{\True $4\%$}
{$0,2\%$}
{$9\%$}
{$2\%$}
\loigiai{
$\delta=|F_{tn}-F_{lt}|/F_{lt}\cdot100\%=4\%$.
}
\end{ex}

% ===== Câu 118 =====
% ID: AIQ_L10C3_FULL_1576
% Mức: VDC
\begin{ex}
% ID: AIQ_L10C3_FULL_1576
% Mức: VDC
Giá trị thực nghiệm là $4,9\ \mathrm{N}$, giá trị lí thuyết là $5,1\ \mathrm{N}$. Sai số tương đối gần nhất bằng bao nhiêu?
\choice
{$0,2\%$}
{\True $3,92\%$}
{$8,92\%$}
{$1,96\%$}
\loigiai{
$\delta=|F_{tn}-F_{lt}|/F_{lt}\cdot100\%=3,92\%$.
}
\end{ex}

% ===== Câu 119 =====
% ID: AIQ_L10C3_FULL_1577
% Mức: NB
\begin{ex}
% ID: AIQ_L10C3_FULL_1577
% Mức: NB
Phân tích sai số trong lượt thí nghiệm 1.
\choiceTF
{\True Ma sát ở ròng rọc có thể làm số chỉ lực kế sai lệch.}
{\True Nhìn xiên mặt chia độ gây sai số thị sai.}
{Đo lặp lại loại bỏ hoàn toàn mọi sai số hệ thống.}
{\True Chỉnh số 0 trước khi đo giúp hạn chế sai số hệ thống.}
\loigiai{
\begin{itemchoice}
\item (Đúng) Ma sát ở ròng rọc có thể làm số chỉ lực kế sai lệch. (Vì): ma sát làm thay đổi lực căng truyền qua ròng rọc. b) Đúng: nhìn xiên gây thị sai. c) Sai: đo lặp không loại bỏ hoàn toàn sai số hệ thống. d) Đúng: chỉnh số 0 giúp hạn chế sai số hệ thống
\item (Đúng) Nhìn xiên mặt chia độ gây sai số thị sai.
\item (Sai) Đo lặp lại loại bỏ hoàn toàn mọi sai số hệ thống.
\item (Đúng) Chỉnh số 0 trước khi đo giúp hạn chế sai số hệ thống.
\end{itemchoice}
}
\end{ex}

% ===== Câu 120 =====
% ID: AIQ_L10C3_FULL_1578
% Mức: TH
\begin{ex}
% ID: AIQ_L10C3_FULL_1578
% Mức: TH
Phân tích sai số trong lượt thí nghiệm 2.
\choiceTF
{\True Ma sát ở ròng rọc có thể làm số chỉ lực kế sai lệch.}
{\True Nhìn xiên mặt chia độ gây sai số thị sai.}
{Đo lặp lại loại bỏ hoàn toàn mọi sai số hệ thống.}
{\True Chỉnh số 0 trước khi đo giúp hạn chế sai số hệ thống.}
\loigiai{
\begin{itemchoice}
\item (Đúng) Ma sát ở ròng rọc có thể làm số chỉ lực kế sai lệch. (Vì): ma sát làm thay đổi lực căng truyền qua ròng rọc. b) Đúng: nhìn xiên gây thị sai. c) Sai: đo lặp không loại bỏ hoàn toàn sai số hệ thống. d) Đúng: chỉnh số 0 giúp hạn chế sai số hệ thống
\item (Đúng) Nhìn xiên mặt chia độ gây sai số thị sai.
\item (Sai) Đo lặp lại loại bỏ hoàn toàn mọi sai số hệ thống.
\item (Đúng) Chỉnh số 0 trước khi đo giúp hạn chế sai số hệ thống.
\end{itemchoice}
}
\end{ex}

% ===== Câu 121 =====
% ID: AIQ_L10C3_FULL_1579
% Mức: VD
\begin{ex}
% ID: AIQ_L10C3_FULL_1579
% Mức: VD
Phân tích sai số trong lượt thí nghiệm 3.
\choiceTF
{\True Ma sát ở ròng rọc có thể làm số chỉ lực kế sai lệch.}
{\True Nhìn xiên mặt chia độ gây sai số thị sai.}
{Đo lặp lại loại bỏ hoàn toàn mọi sai số hệ thống.}
{\True Chỉnh số 0 trước khi đo giúp hạn chế sai số hệ thống.}
\loigiai{
\begin{itemchoice}
\item (Đúng) Ma sát ở ròng rọc có thể làm số chỉ lực kế sai lệch. (Vì): ma sát làm thay đổi lực căng truyền qua ròng rọc. b) Đúng: nhìn xiên gây thị sai. c) Sai: đo lặp không loại bỏ hoàn toàn sai số hệ thống. d) Đúng: chỉnh số 0 giúp hạn chế sai số hệ thống
\item (Đúng) Nhìn xiên mặt chia độ gây sai số thị sai.
\item (Sai) Đo lặp lại loại bỏ hoàn toàn mọi sai số hệ thống.
\item (Đúng) Chỉnh số 0 trước khi đo giúp hạn chế sai số hệ thống.
\end{itemchoice}
}
\end{ex}

% ===== Câu 122 =====
% ID: AIQ_L10C3_FULL_1580
% Mức: VD
\begin{ex}
% ID: AIQ_L10C3_FULL_1580
% Mức: VD
Phân tích sai số trong lượt thí nghiệm 4.
\choiceTF
{\True Ma sát ở ròng rọc có thể làm số chỉ lực kế sai lệch.}
{\True Nhìn xiên mặt chia độ gây sai số thị sai.}
{Đo lặp lại loại bỏ hoàn toàn mọi sai số hệ thống.}
{\True Chỉnh số 0 trước khi đo giúp hạn chế sai số hệ thống.}
\loigiai{
\begin{itemchoice}
\item (Đúng) Ma sát ở ròng rọc có thể làm số chỉ lực kế sai lệch. (Vì): ma sát làm thay đổi lực căng truyền qua ròng rọc. b) Đúng: nhìn xiên gây thị sai. c) Sai: đo lặp không loại bỏ hoàn toàn sai số hệ thống. d) Đúng: chỉnh số 0 giúp hạn chế sai số hệ thống
\item (Đúng) Nhìn xiên mặt chia độ gây sai số thị sai.
\item (Sai) Đo lặp lại loại bỏ hoàn toàn mọi sai số hệ thống.
\item (Đúng) Chỉnh số 0 trước khi đo giúp hạn chế sai số hệ thống.
\end{itemchoice}
}
\end{ex}

% ===== Câu 123 =====
% ID: AIQ_L10C3_FULL_1581
% Mức: VDC
\begin{ex}
% ID: AIQ_L10C3_FULL_1581
% Mức: VDC
Phân tích sai số trong lượt thí nghiệm 5.
\choiceTF
{\True Ma sát ở ròng rọc có thể làm số chỉ lực kế sai lệch.}
{\True Nhìn xiên mặt chia độ gây sai số thị sai.}
{Đo lặp lại loại bỏ hoàn toàn mọi sai số hệ thống.}
{\True Chỉnh số 0 trước khi đo giúp hạn chế sai số hệ thống.}
\loigiai{
\begin{itemchoice}
\item (Đúng) Ma sát ở ròng rọc có thể làm số chỉ lực kế sai lệch. (Vì): ma sát làm thay đổi lực căng truyền qua ròng rọc. b) Đúng: nhìn xiên gây thị sai. c) Sai: đo lặp không loại bỏ hoàn toàn sai số hệ thống. d) Đúng: chỉnh số 0 giúp hạn chế sai số hệ thống
\item (Đúng) Nhìn xiên mặt chia độ gây sai số thị sai.
\item (Sai) Đo lặp lại loại bỏ hoàn toàn mọi sai số hệ thống.
\item (Đúng) Chỉnh số 0 trước khi đo giúp hạn chế sai số hệ thống.
\end{itemchoice}
}
\end{ex}

% ===== Câu 124 =====
% ID: AIQ_L10C3_FULL_1582
% Mức: TH
\begin{ex}
% ID: AIQ_L10C3_FULL_1582
% Mức: TH
Hợp lực lí thuyết $5N$, đo được $4,8$ N. Tính sai số tương đối theo %, làm tròn đến hai chữ số thập phân.
\shortans{4}
\loigiai{
Sai số tương đối được tính bởi $\delta=|F_{tn}-F_{lt}|/F_{lt}\cdot100\%=4\%$. Đáp án không ghi đơn vị.
}
\end{ex}

% ===== Câu 125 =====
% ID: AIQ_L10C3_FULL_1583
% Mức: TH
\begin{ex}
% ID: AIQ_L10C3_FULL_1583
% Mức: TH
Hợp lực lí thuyết $6N$, đo được $5,8$ N. Tính sai số tương đối theo %, làm tròn đến hai chữ số thập phân.
\shortans{3,33}
\loigiai{
Sai số tương đối được tính bởi $\delta=|F_{tn}-F_{lt}|/F_{lt}\cdot100\%=3,33\%$. Đáp án không ghi đơn vị.
}
\end{ex}

% ===== Câu 126 =====
% ID: AIQ_L10C3_FULL_1584
% Mức: VD
\begin{ex}
% ID: AIQ_L10C3_FULL_1584
% Mức: VD
Hợp lực lí thuyết $7N$, đo được $6,8$ N. Tính sai số tương đối theo %, làm tròn đến hai chữ số thập phân.
\shortans{2,86}
\loigiai{
Sai số tương đối được tính bởi $\delta=|F_{tn}-F_{lt}|/F_{lt}\cdot100\%=2,86\%$. Đáp án không ghi đơn vị.
}
\end{ex}

% ===== Câu 127 =====
% ID: AIQ_L10C3_FULL_1585
% Mức: VD
\begin{ex}
% ID: AIQ_L10C3_FULL_1585
% Mức: VD
Hợp lực lí thuyết $8N$, đo được $7,8$ N. Tính sai số tương đối theo %, làm tròn đến hai chữ số thập phân.
\shortans{2,5}
\loigiai{
Sai số tương đối được tính bởi $\delta=|F_{tn}-F_{lt}|/F_{lt}\cdot100\%=2,5\%$. Đáp án không ghi đơn vị.
}
\end{ex}

% ===== Câu 128 =====
% ID: AIQ_L10C3_FULL_1586
% Mức: VD
\begin{ex}
% ID: AIQ_L10C3_FULL_1586
% Mức: VD
Hợp lực lí thuyết $9N$, đo được $8,8$ N. Tính sai số tương đối theo %, làm tròn đến hai chữ số thập phân.
\shortans{2,22}
\loigiai{
Sai số tương đối được tính bởi $\delta=|F_{tn}-F_{lt}|/F_{lt}\cdot100\%=2,22\%$. Đáp án không ghi đơn vị.
}
\end{ex}

% ===== Câu 129 =====
% ID: AIQ_L10C3_FULL_1587
% Mức: VDC
\begin{ex}
% ID: AIQ_L10C3_FULL_1587
% Mức: VDC
Hợp lực lí thuyết $10N$, đo được $9,8$ N. Tính sai số tương đối theo %, làm tròn đến hai chữ số thập phân.
\shortans{2}
\loigiai{
Sai số tương đối được tính bởi $\delta=|F_{tn}-F_{lt}|/F_{lt}\cdot100\%=2\%$. Đáp án không ghi đơn vị.
}
\end{ex}

% ===== Câu 130 =====
% ID: AIQ_L10C3_FULL_1588
% Mức: TH
\begin{bt}
% ID: AIQ_L10C3_FULL_1588
% Mức: TH
Giá trị lí thuyết $6N$, thực nghiệm $5,7$ N. Tính sai số tương đối và nêu hai nguyên nhân có thể gây sai lệch.
\loigiai{
Dùng $\delta=|F_{tn}-F_{lt}|/F_{lt}\cdot100\%$ thu được 5%. Sai lệch có thể do ma sát ở ròng rọc, thị sai, dây không đồng phẳng hoặc lực kế chưa chỉnh số 0.
}
\end{bt}

% ===== Câu 131 =====
% ID: AIQ_L10C3_FULL_1589
% Mức: TH
\begin{bt}
% ID: AIQ_L10C3_FULL_1589
% Mức: TH
Giá trị lí thuyết $7N$, thực nghiệm $6,7$ N. Tính sai số tương đối và nêu hai nguyên nhân có thể gây sai lệch.
\loigiai{
Dùng $\delta=|F_{tn}-F_{lt}|/F_{lt}\cdot100\%$ thu được 4,29%. Sai lệch có thể do ma sát ở ròng rọc, thị sai, dây không đồng phẳng hoặc lực kế chưa chỉnh số 0.
}
\end{bt}

% ===== Câu 132 =====
% ID: AIQ_L10C3_FULL_1590
% Mức: VD
\begin{bt}
% ID: AIQ_L10C3_FULL_1590
% Mức: VD
Giá trị lí thuyết $8N$, thực nghiệm $7,7$ N. Tính sai số tương đối và nêu hai nguyên nhân có thể gây sai lệch.
\loigiai{
Dùng $\delta=|F_{tn}-F_{lt}|/F_{lt}\cdot100\%$ thu được 3,75%. Sai lệch có thể do ma sát ở ròng rọc, thị sai, dây không đồng phẳng hoặc lực kế chưa chỉnh số 0.
}
\end{bt}

% ===== Câu 133 =====
% ID: AIQ_L10C3_FULL_1591
% Mức: VD
\begin{bt}
% ID: AIQ_L10C3_FULL_1591
% Mức: VD
Giá trị lí thuyết $9N$, thực nghiệm $8,7$ N. Tính sai số tương đối và nêu hai nguyên nhân có thể gây sai lệch.
\loigiai{
Dùng $\delta=|F_{tn}-F_{lt}|/F_{lt}\cdot100\%$ thu được 3,33%. Sai lệch có thể do ma sát ở ròng rọc, thị sai, dây không đồng phẳng hoặc lực kế chưa chỉnh số 0.
}
\end{bt}

% ===== Câu 134 =====
% ID: AIQ_L10C3_FULL_1592
% Mức: VD
\begin{bt}
% ID: AIQ_L10C3_FULL_1592
% Mức: VD
Giá trị lí thuyết $10N$, thực nghiệm $9,7$ N. Tính sai số tương đối và nêu hai nguyên nhân có thể gây sai lệch.
\loigiai{
Dùng $\delta=|F_{tn}-F_{lt}|/F_{lt}\cdot100\%$ thu được 3%. Sai lệch có thể do ma sát ở ròng rọc, thị sai, dây không đồng phẳng hoặc lực kế chưa chỉnh số 0.
}
\end{bt}

% ===== Câu 135 =====
% ID: AIQ_L10C3_FULL_1593
% Mức: VDC
\begin{bt}
% ID: AIQ_L10C3_FULL_1593
% Mức: VDC
Giá trị lí thuyết $11N$, thực nghiệm $10,7$ N. Tính sai số tương đối và nêu hai nguyên nhân có thể gây sai lệch.
\loigiai{
Dùng $\delta=|F_{tn}-F_{lt}|/F_{lt}\cdot100\%$ thu được 2,73%. Sai lệch có thể do ma sát ở ròng rọc, thị sai, dây không đồng phẳng hoặc lực kế chưa chỉnh số 0.
}
\end{bt}

% ===== Câu 136 =====
% ID: AIQ_L10C3_FULL_1594
% Mức: NB
\dangbt{Thiết kế, cải tiến và báo cáo thí nghiệm tổng hợp lực}
\begin{ex}
% ID: AIQ_L10C3_FULL_1594
% Mức: NB
Trong báo cáo thí nghiệm, phần nào nêu mục tiêu cần kiểm chứng?
\choice
{\True Mục đích thí nghiệm}
{Kết luận cá nhân không căn cứ}
{Danh sách màu dây}
{Trang trí bìa}
\loigiai{
Mục đích nêu rõ đại lượng hoặc quy tắc cần kiểm chứng.
}
\end{ex}

% ===== Câu 137 =====
% ID: AIQ_L10C3_FULL_1595
% Mức: NB
\begin{ex}
% ID: AIQ_L10C3_FULL_1595
% Mức: NB
Một phương án tốt cần kiểm soát yếu tố nào?
\choice
{Màu áo người đo}
{\True Vị trí vòng nhẫn và góc giữa các lực}
{Thời gian trong ngày}
{Tên phòng học}
\loigiai{
Các yếu tố trực tiếp ảnh hưởng phép đo phải được kiểm soát.
}
\end{ex}

% ===== Câu 138 =====
% ID: AIQ_L10C3_FULL_1596
% Mức: NB
\begin{ex}
% ID: AIQ_L10C3_FULL_1596
% Mức: NB
Để tăng độ tin cậy của kết quả, nên làm gì?
\choice
{Chỉ đo một lần}
{Bỏ qua số liệu lệch}
{\True Lặp lại phép đo và lấy trung bình}
{Đổi kết quả theo dự đoán}
\loigiai{
Đo lặp và xử lí thống kê làm kết quả đáng tin cậy hơn.
}
\end{ex}

% ===== Câu 139 =====
% ID: AIQ_L10C3_FULL_1597
% Mức: TH
\begin{ex}
% ID: AIQ_L10C3_FULL_1597
% Mức: TH
Cải tiến nào giúp giảm ma sát tại điểm đổi hướng dây?
\choice
{Buộc nút dây chặt hơn vào bảng}
{Tăng khối lượng vòng nhẫn}
{Dùng dây rất nhám}
{\True Dùng ròng rọc quay nhẹ}
\loigiai{
Ròng rọc quay nhẹ làm giảm ma sát cản trở.
}
\end{ex}

% ===== Câu 140 =====
% ID: AIQ_L10C3_FULL_1598
% Mức: TH
\begin{ex}
% ID: AIQ_L10C3_FULL_1598
% Mức: TH
Kết luận khoa học phải dựa chủ yếu vào đâu?
\choice
{\True Số liệu và sai số đo}
{Cảm giác người làm}
{Màu biểu đồ}
{Dự đoán ban đầu}
\loigiai{
Kết luận phải phù hợp dữ liệu và độ không đảm bảo.
}
\end{ex}

% ===== Câu 141 =====
% ID: AIQ_L10C3_FULL_1599
% Mức: TH
\begin{ex}
% ID: AIQ_L10C3_FULL_1599
% Mức: TH
Biểu đồ vectơ trong báo cáo cần thể hiện gì?
\choice
{Chỉ màu sắc}
{\True Tỉ xích, phương, chiều và độ lớn}
{Chỉ tên lực}
{Chỉ số thứ tự}
\loigiai{
Bốn yếu tố này giúp biểu diễn vectơ có nghĩa.
}
\end{ex}

% ===== Câu 142 =====
% ID: AIQ_L10C3_FULL_1600
% Mức: VD
\begin{ex}
% ID: AIQ_L10C3_FULL_1600
% Mức: VD
Khi thay đổi góc giữa hai lực, cần giữ đại lượng nào nếu muốn khảo sát riêng ảnh hưởng của góc?
\choice
{Màu dây}
{Tên lực kế}
{\True Độ lớn hai lực}
{Kích thước giấy}
\loigiai{
Giữ độ lớn lực để góc là biến độc lập.
}
\end{ex}

% ===== Câu 143 =====
% ID: AIQ_L10C3_FULL_1601
% Mức: VD
\begin{ex}
% ID: AIQ_L10C3_FULL_1601
% Mức: VD
Trình tự hợp lí là gì?
\choice
{Kết luận, đo, bố trí, xử lí}
{Đo, kết luận, bố trí, xử lí}
{Xử lí, bố trí, kết luận, đo}
{\True Bố trí, đo, xử lí, kết luận}
\loigiai{
Tiến trình thực nghiệm bắt đầu từ bố trí và kết thúc bằng kết luận.
}
\end{ex}

% ===== Câu 144 =====
% ID: AIQ_L10C3_FULL_1602
% Mức: VD
\begin{ex}
% ID: AIQ_L10C3_FULL_1602
% Mức: VD
Ảnh chụp bố trí thí nghiệm trong báo cáo có tác dụng gì?
\choice
{\True Minh chứng cách lắp đặt và vị trí dụng cụ}
{Thay toàn bộ số liệu}
{Thay lời giải thích sai số}
{Làm tăng lực đo}
\loigiai{
Ảnh hỗ trợ kiểm chứng bố trí nhưng không thay thế dữ liệu.
}
\end{ex}

% ===== Câu 145 =====
% ID: AIQ_L10C3_FULL_1603
% Mức: VDC
\begin{ex}
% ID: AIQ_L10C3_FULL_1603
% Mức: VDC
Khi kết quả lệch dự đoán, cách xử lí đúng là gì?
\choice
{Sửa số liệu cho khớp}
{\True Kiểm tra bố trí, đo lại và phân tích sai số}
{Xóa toàn bộ báo cáo}
{Giữ kết quả mà không nhận xét}
\loigiai{
Cần tìm nguyên nhân và đo kiểm lại một cách minh bạch.
}
\end{ex}

% ===== Câu 146 =====
% ID: AIQ_L10C3_FULL_1604
% Mức: NB
\begin{ex}
% ID: AIQ_L10C3_FULL_1604
% Mức: NB
Đánh giá phương án và báo cáo thí nghiệm 1.
\choiceTF
{\True Báo cáo cần nêu mục tiêu, dụng cụ, cách tiến hành và kết quả.}
{Có thể thay số liệu đo bằng số liệu dự đoán.}
{\True Nên ghi rõ tỉ xích khi vẽ vectơ.}
{\True Kết luận cần đối chiếu kết quả thực nghiệm với lí thuyết.}
\loigiai{
\begin{itemchoice}
\item (Đúng) Báo cáo cần nêu mục tiêu, dụng cụ, cách tiến hành và kết quả. (Vì): báo cáo phải đủ mục tiêu, dụng cụ, cách tiến hành và kết quả. b) Sai: không được thay số liệu đo bằng dự đoán. c) Đúng: tỉ xích là thông tin bắt buộc khi vẽ vectơ. d) Đúng: kết luận phải dựa trên đối chiếu thực nghiệm với lí thuyết
\item (Sai) Có thể thay số liệu đo bằng số liệu dự đoán.
\item (Đúng) Nên ghi rõ tỉ xích khi vẽ vectơ.
\item (Đúng) Kết luận cần đối chiếu kết quả thực nghiệm với lí thuyết.
\end{itemchoice}
}
\end{ex}

% ===== Câu 147 =====
% ID: AIQ_L10C3_FULL_1605
% Mức: TH
\begin{ex}
% ID: AIQ_L10C3_FULL_1605
% Mức: TH
Đánh giá phương án và báo cáo thí nghiệm 2.
\choiceTF
{\True Báo cáo cần nêu mục tiêu, dụng cụ, cách tiến hành và kết quả.}
{Có thể thay số liệu đo bằng số liệu dự đoán.}
{\True Nên ghi rõ tỉ xích khi vẽ vectơ.}
{\True Kết luận cần đối chiếu kết quả thực nghiệm với lí thuyết.}
\loigiai{
\begin{itemchoice}
\item (Đúng) Báo cáo cần nêu mục tiêu, dụng cụ, cách tiến hành và kết quả. (Vì): báo cáo phải đủ mục tiêu, dụng cụ, cách tiến hành và kết quả. b) Sai: không được thay số liệu đo bằng dự đoán. c) Đúng: tỉ xích là thông tin bắt buộc khi vẽ vectơ. d) Đúng: kết luận phải dựa trên đối chiếu thực nghiệm với lí thuyết
\item (Sai) Có thể thay số liệu đo bằng số liệu dự đoán.
\item (Đúng) Nên ghi rõ tỉ xích khi vẽ vectơ.
\item (Đúng) Kết luận cần đối chiếu kết quả thực nghiệm với lí thuyết.
\end{itemchoice}
}
\end{ex}

% ===== Câu 148 =====
% ID: AIQ_L10C3_FULL_1606
% Mức: VD
\begin{ex}
% ID: AIQ_L10C3_FULL_1606
% Mức: VD
Đánh giá phương án và báo cáo thí nghiệm 3.
\choiceTF
{\True Báo cáo cần nêu mục tiêu, dụng cụ, cách tiến hành và kết quả.}
{Có thể thay số liệu đo bằng số liệu dự đoán.}
{\True Nên ghi rõ tỉ xích khi vẽ vectơ.}
{\True Kết luận cần đối chiếu kết quả thực nghiệm với lí thuyết.}
\loigiai{
\begin{itemchoice}
\item (Đúng) Báo cáo cần nêu mục tiêu, dụng cụ, cách tiến hành và kết quả. (Vì): báo cáo phải đủ mục tiêu, dụng cụ, cách tiến hành và kết quả. b) Sai: không được thay số liệu đo bằng dự đoán. c) Đúng: tỉ xích là thông tin bắt buộc khi vẽ vectơ. d) Đúng: kết luận phải dựa trên đối chiếu thực nghiệm với lí thuyết
\item (Sai) Có thể thay số liệu đo bằng số liệu dự đoán.
\item (Đúng) Nên ghi rõ tỉ xích khi vẽ vectơ.
\item (Đúng) Kết luận cần đối chiếu kết quả thực nghiệm với lí thuyết.
\end{itemchoice}
}
\end{ex}

% ===== Câu 149 =====
% ID: AIQ_L10C3_FULL_1607
% Mức: VD
\begin{ex}
% ID: AIQ_L10C3_FULL_1607
% Mức: VD
Đánh giá phương án và báo cáo thí nghiệm 4.
\choiceTF
{\True Báo cáo cần nêu mục tiêu, dụng cụ, cách tiến hành và kết quả.}
{Có thể thay số liệu đo bằng số liệu dự đoán.}
{\True Nên ghi rõ tỉ xích khi vẽ vectơ.}
{\True Kết luận cần đối chiếu kết quả thực nghiệm với lí thuyết.}
\loigiai{
\begin{itemchoice}
\item (Đúng) Báo cáo cần nêu mục tiêu, dụng cụ, cách tiến hành và kết quả. (Vì): báo cáo phải đủ mục tiêu, dụng cụ, cách tiến hành và kết quả. b) Sai: không được thay số liệu đo bằng dự đoán. c) Đúng: tỉ xích là thông tin bắt buộc khi vẽ vectơ. d) Đúng: kết luận phải dựa trên đối chiếu thực nghiệm với lí thuyết
\item (Sai) Có thể thay số liệu đo bằng số liệu dự đoán.
\item (Đúng) Nên ghi rõ tỉ xích khi vẽ vectơ.
\item (Đúng) Kết luận cần đối chiếu kết quả thực nghiệm với lí thuyết.
\end{itemchoice}
}
\end{ex}

% ===== Câu 150 =====
% ID: AIQ_L10C3_FULL_1608
% Mức: VDC
\begin{ex}
% ID: AIQ_L10C3_FULL_1608
% Mức: VDC
Đánh giá phương án và báo cáo thí nghiệm 5.
\choiceTF
{\True Báo cáo cần nêu mục tiêu, dụng cụ, cách tiến hành và kết quả.}
{Có thể thay số liệu đo bằng số liệu dự đoán.}
{\True Nên ghi rõ tỉ xích khi vẽ vectơ.}
{\True Kết luận cần đối chiếu kết quả thực nghiệm với lí thuyết.}
\loigiai{
\begin{itemchoice}
\item (Đúng) Báo cáo cần nêu mục tiêu, dụng cụ, cách tiến hành và kết quả. (Vì): báo cáo phải đủ mục tiêu, dụng cụ, cách tiến hành và kết quả. b) Sai: không được thay số liệu đo bằng dự đoán. c) Đúng: tỉ xích là thông tin bắt buộc khi vẽ vectơ. d) Đúng: kết luận phải dựa trên đối chiếu thực nghiệm với lí thuyết
\item (Sai) Có thể thay số liệu đo bằng số liệu dự đoán.
\item (Đúng) Nên ghi rõ tỉ xích khi vẽ vectơ.
\item (Đúng) Kết luận cần đối chiếu kết quả thực nghiệm với lí thuyết.
\end{itemchoice}
}
\end{ex}

% ===== Câu 151 =====
% ID: AIQ_L10C3_FULL_1609
% Mức: TH
\begin{ex}
% ID: AIQ_L10C3_FULL_1609
% Mức: TH
Một phương án đo lặp 3 lần cho mỗi cấu hình và khảo sát 4 cấu hình góc. Tổng số lượt đo là bao nhiêu?
\shortans{12}
\loigiai{
Mỗi cấu hình được đo 3 lần và có 4 cấu hình, nên tổng số lượt đo bằng tích hai đại lượng, kết quả 12. Đáp án không ghi đơn vị.
}
\end{ex}

% ===== Câu 152 =====
% ID: AIQ_L10C3_FULL_1610
% Mức: TH
\begin{ex}
% ID: AIQ_L10C3_FULL_1610
% Mức: TH
Một phương án đo lặp 4 lần cho mỗi cấu hình và khảo sát 4 cấu hình góc. Tổng số lượt đo là bao nhiêu?
\shortans{16}
\loigiai{
Mỗi cấu hình được đo 4 lần và có 4 cấu hình, nên tổng số lượt đo bằng tích hai đại lượng, kết quả 16. Đáp án không ghi đơn vị.
}
\end{ex}

% ===== Câu 153 =====
% ID: AIQ_L10C3_FULL_1611
% Mức: VD
\begin{ex}
% ID: AIQ_L10C3_FULL_1611
% Mức: VD
Một phương án đo lặp 5 lần cho mỗi cấu hình và khảo sát 4 cấu hình góc. Tổng số lượt đo là bao nhiêu?
\shortans{20}
\loigiai{
Mỗi cấu hình được đo 5 lần và có 4 cấu hình, nên tổng số lượt đo bằng tích hai đại lượng, kết quả 20. Đáp án không ghi đơn vị.
}
\end{ex}

% ===== Câu 154 =====
% ID: AIQ_L10C3_FULL_1612
% Mức: VD
\begin{ex}
% ID: AIQ_L10C3_FULL_1612
% Mức: VD
Một phương án đo lặp 6 lần cho mỗi cấu hình và khảo sát 4 cấu hình góc. Tổng số lượt đo là bao nhiêu?
\shortans{24}
\loigiai{
Mỗi cấu hình được đo 6 lần và có 4 cấu hình, nên tổng số lượt đo bằng tích hai đại lượng, kết quả 24. Đáp án không ghi đơn vị.
}
\end{ex}

% ===== Câu 155 =====
% ID: AIQ_L10C3_FULL_1613
% Mức: VD
\begin{ex}
% ID: AIQ_L10C3_FULL_1613
% Mức: VD
Một phương án đo lặp 7 lần cho mỗi cấu hình và khảo sát 4 cấu hình góc. Tổng số lượt đo là bao nhiêu?
\shortans{28}
\loigiai{
Mỗi cấu hình được đo 7 lần và có 4 cấu hình, nên tổng số lượt đo bằng tích hai đại lượng, kết quả 28. Đáp án không ghi đơn vị.
}
\end{ex}

% ===== Câu 156 =====
% ID: AIQ_L10C3_FULL_1614
% Mức: VDC
\begin{ex}
% ID: AIQ_L10C3_FULL_1614
% Mức: VDC
Một phương án đo lặp 8 lần cho mỗi cấu hình và khảo sát 4 cấu hình góc. Tổng số lượt đo là bao nhiêu?
\shortans{32}
\loigiai{
Mỗi cấu hình được đo 8 lần và có 4 cấu hình, nên tổng số lượt đo bằng tích hai đại lượng, kết quả 32. Đáp án không ghi đơn vị.
}
\end{ex}

% ===== Câu 157 =====
% ID: AIQ_L10C3_FULL_1615
% Mức: TH
\begin{ex}
% ID: AIQ_L10C3_FULL_1615
% Mức: TH
Thiết kế quy trình kiểm chứng quy tắc hình bình hành với ít nhất ba giá trị góc khác nhau (phương án 1).
\choiceTF

\loigiai{
Giữ độ lớn hai lực không đổi; lần lượt đặt ít nhất ba góc; mỗi góc đo lực cân bằng nhiều lần và lấy trung bình; dựng hợp lực lí thuyết; tính sai số tương đối; lập bảng, biểu đồ và kết luận trong giới hạn sai số.
}
\end{ex}

% ===== Câu 158 =====
% ID: AIQ_L10C3_FULL_1616
% Mức: TH
\begin{ex}
% ID: AIQ_L10C3_FULL_1616
% Mức: TH
Thiết kế quy trình kiểm chứng quy tắc hình bình hành với ít nhất ba giá trị góc khác nhau (phương án 2).
\choiceTF

\loigiai{
Giữ độ lớn hai lực không đổi; lần lượt đặt ít nhất ba góc; mỗi góc đo lực cân bằng nhiều lần và lấy trung bình; dựng hợp lực lí thuyết; tính sai số tương đối; lập bảng, biểu đồ và kết luận trong giới hạn sai số.
}
\end{ex}

% ===== Câu 159 =====
% ID: AIQ_L10C3_FULL_1617
% Mức: VD
\begin{ex}
% ID: AIQ_L10C3_FULL_1617
% Mức: VD
Thiết kế quy trình kiểm chứng quy tắc hình bình hành với ít nhất ba giá trị góc khác nhau (phương án 3).
\choiceTF

\loigiai{
Giữ độ lớn hai lực không đổi; lần lượt đặt ít nhất ba góc; mỗi góc đo lực cân bằng nhiều lần và lấy trung bình; dựng hợp lực lí thuyết; tính sai số tương đối; lập bảng, biểu đồ và kết luận trong giới hạn sai số.
}
\end{ex}

% ===== Câu 160 =====
% ID: AIQ_L10C3_FULL_1618
% Mức: VD
\begin{ex}
% ID: AIQ_L10C3_FULL_1618
% Mức: VD
Thiết kế quy trình kiểm chứng quy tắc hình bình hành với ít nhất ba giá trị góc khác nhau (phương án 4).
\choiceTF

\loigiai{
Giữ độ lớn hai lực không đổi; lần lượt đặt ít nhất ba góc; mỗi góc đo lực cân bằng nhiều lần và lấy trung bình; dựng hợp lực lí thuyết; tính sai số tương đối; lập bảng, biểu đồ và kết luận trong giới hạn sai số.
}
\end{ex}

% ===== Câu 161 =====
% ID: AIQ_L10C3_FULL_1619
% Mức: VD
\begin{ex}
% ID: AIQ_L10C3_FULL_1619
% Mức: VD
Thiết kế quy trình kiểm chứng quy tắc hình bình hành với ít nhất ba giá trị góc khác nhau (phương án 5).
\choiceTF

\loigiai{
Giữ độ lớn hai lực không đổi; lần lượt đặt ít nhất ba góc; mỗi góc đo lực cân bằng nhiều lần và lấy trung bình; dựng hợp lực lí thuyết; tính sai số tương đối; lập bảng, biểu đồ và kết luận trong giới hạn sai số.
}
\end{ex}

% ===== Câu 162 =====
% ID: AIQ_L10C3_FULL_1620
% Mức: VDC
\begin{ex}
% ID: AIQ_L10C3_FULL_1620
% Mức: VDC
Thiết kế quy trình kiểm chứng quy tắc hình bình hành với ít nhất ba giá trị góc khác nhau (phương án 6).
\choiceTF

\loigiai{
Giữ độ lớn hai lực không đổi; lần lượt đặt ít nhất ba góc; mỗi góc đo lực cân bằng nhiều lần và lấy trung bình; dựng hợp lực lí thuyết; tính sai số tương đối; lập bảng, biểu đồ và kết luận trong giới hạn sai số.
}
\end{ex}
